% =====================================================================
%  Anatomía de un LLM — Fase VI — Módulo 28
%  LLM como optimizadores y del prototipo a la producción
%  Compilar: latexmk -r ../latexmkrc M28-llm-produccion.tex
% =====================================================================
\documentclass[11pt,a4paper,twoside,openany]{book}
\usepackage{../estilo-apuntes}

\newcommand{\ApuntesCurso}{Anatomía de un LLM}
\newcommand{\ApuntesModulo}{Módulo 28}

% Notación propia del módulo
\DeclareMathOperator*{\argmin}{arg\,min}
\DeclareMathOperator*{\argmax}{arg\,max}
\DeclareMathOperator{\VAN}{VAN}
\DeclareMathOperator{\PSI}{PSI}
\DeclareMathOperator{\KL}{KL}
\DeclareMathOperator{\Var}{Var}
\newcommand{\T}{^{\mathsf{T}}}
\newcommand{\Prob}{\mathbb{P}}
\newcommand{\Ind}{\mathbf{1}}
\newcommand{\eur}{\text{\,€}}
\newcommand{\Prog}{\mathcal{P}}
\makeatletter\g@addto@macro\UrlBreaks{\do\-\do\/}\makeatother

\hypersetup{pdftitle={Módulo 28. LLM como optimizadores y del prototipo a la producción (resumen)},pdfauthor={Anatomía de un LLM}}

\makeatletter
\newlabel{fig:funsearch}{{1.1}{2}{Bucle de FunSearch. El LLM nunca ve el evaluador: solo ve programas anteriores ordenados por su puntuación}{figure.1.1}{}}
\newlabel{fig:funsearch@cref}{{[figure][1][1]1.1}{[1][2][]2}}
\newlabel{ej:ae-resultados}{{1.7}{4}{Resultados publicados de AlphaEvolve}{ejemplo.7}{}}
\newlabel{ej:ae-resultados@cref}{{[ejemplo][7][]1.7}{[1][4][]4}}
\newlabel{ej:ff-bf}{{1.12}{5}{FF frente a BF a mano}{ejemplo.12}{}}
\newlabel{ej:ff-bf@cref}{{[ejemplo][12][]1.12}{[1][5][]5}}
\newlabel{ej:caso-bp}{{1.13}{6}{Búsqueda evolutiva de una regla de prioridad}{ejemplo.13}{}}
\newlabel{ej:caso-bp@cref}{{[ejemplo][13][]1.13}{[1][5][]6}}
\newlabel{obs:errores1}{{1.14}{7}{Errores frecuentes en la práctica}{observacion.14}{}}
\newlabel{obs:errores1@cref}{{[observacion][14][]1.14}{[1][7][]7}}
\newlabel{eje:1-componentes}{{1.1}{8}{}{ejercicio.1.1}{}}
\newlabel{eje:1-componentes@cref}{{[ejercicio][1][1]1.1}{[1][7][]8}}
\newlabel{eje:1-softmax}{{1.2}{8}{}{ejercicio.1.2}{}}
\newlabel{eje:1-softmax@cref}{{[ejercicio][2][1]1.2}{[1][8][]8}}
\newlabel{eje:1-ffbf}{{1.3}{8}{}{ejercicio.1.3}{}}
\newlabel{eje:1-ffbf@cref}{{[ejercicio][3][1]1.3}{[1][8][]8}}
\newlabel{eje:1-nextfit}{{1.4}{8}{}{ejercicio.1.4}{}}
\newlabel{eje:1-nextfit@cref}{{[ejercicio][4][1]1.4}{[1][8][]8}}
\newlabel{eje:1-presupuesto}{{1.5}{8}{}{ejercicio.1.5}{}}
\newlabel{eje:1-presupuesto@cref}{{[ejercicio][5][1]1.5}{[1][8][]8}}
\newlabel{eje:1-gramatica}{{1.6}{8}{}{ejercicio.1.6}{}}
\newlabel{eje:1-gramatica@cref}{{[ejercicio][6][1]1.6}{[1][8][]8}}
\newlabel{eje:1-interpretar}{{1.7}{8}{}{ejercicio.1.7}{}}
\newlabel{eje:1-interpretar@cref}{{[ejercicio][7][1]1.7}{[1][8][]8}}
\newlabel{eje:1-diseno}{{1.8}{8}{}{ejercicio.1.8}{}}
\newlabel{eje:1-diseno@cref}{{[ejercicio][8][1]1.8}{[1][8][]8}}
\newlabel{ej:maldicion}{{2.5}{11}{Veinte instrucciones igual de buenas}{ejemplo.5}{}}
\newlabel{ej:maldicion@cref}{{[ejemplo][5][]2.5}{[1][11][]11}}
\newlabel{ej:caso-prompt}{{2.7}{12}{Optimizar las instrucciones de un clasificador}{ejemplo.7}{}}
\newlabel{ej:caso-prompt@cref}{{[ejemplo][7][]2.7}{[1][12][]12}}
\newlabel{ej:caso-codigo}{{2.11}{14}{Caso de estudio: acelerar un servicio de planificación}{ejemplo.11}{}}
\newlabel{ej:caso-codigo@cref}{{[ejemplo][11][]2.11}{[1][13][]14}}
\newlabel{obs:errores2}{{2.12}{14}{Errores frecuentes en la práctica}{observacion.12}{}}
\newlabel{obs:errores2@cref}{{[observacion][12][]2.12}{[1][14][]14}}
\newlabel{eje:2-meta}{{2.1}{15}{}{ejercicio.2.1}{}}
\newlabel{eje:2-meta@cref}{{[ejercicio][1][2]2.1}{[1][15][]15}}
\newlabel{eje:2-metrica}{{2.2}{15}{}{ejercicio.2.2}{}}
\newlabel{eje:2-metrica@cref}{{[ejercicio][2][2]2.2}{[1][15][]15}}
\newlabel{eje:2-maldicion}{{2.3}{15}{}{ejercicio.2.3}{}}
\newlabel{eje:2-maldicion@cref}{{[ejercicio][3][2]2.3}{[1][15][]15}}
\newlabel{eje:2-voto}{{2.4}{15}{}{ejercicio.2.4}{}}
\newlabel{eje:2-voto@cref}{{[ejercicio][4][2]2.4}{[1][15][]15}}
\newlabel{eje:2-amdahl}{{2.5}{15}{}{ejercicio.2.5}{}}
\newlabel{eje:2-amdahl@cref}{{[ejercicio][5][2]2.5}{[1][15][]15}}
\newlabel{eje:2-medida}{{2.6}{15}{}{ejercicio.2.6}{}}
\newlabel{eje:2-medida@cref}{{[ejercicio][6][2]2.6}{[1][15][]15}}
\newlabel{eje:2-cota}{{2.7}{16}{}{ejercicio.2.7}{}}
\newlabel{eje:2-cota@cref}{{[ejercicio][7][2]2.7}{[1][16][]16}}
\newlabel{eje:2-diseno}{{2.8}{16}{}{ejercicio.2.8}{}}
\newlabel{eje:2-diseno@cref}{{[ejercicio][8][2]2.8}{[1][16][]16}}
\newlabel{ej:caso-automl}{{3.9}{19}{Selección del modelo de riesgo de avería}{ejemplo.9}{}}
\newlabel{ej:caso-automl@cref}{{[ejemplo][9][]3.9}{[1][19][]19}}
\newlabel{obs:errores3}{{3.10}{20}{Errores frecuentes en la práctica}{observacion.10}{}}
\newlabel{obs:errores3@cref}{{[observacion][10][]3.10}{[1][20][]20}}
\newlabel{eje:3-dimensiones}{{3.1}{21}{}{ejercicio.3.1}{}}
\newlabel{eje:3-dimensiones@cref}{{[ejercicio][1][3]3.1}{[1][21][]21}}
\newlabel{eje:3-rejilla}{{3.2}{21}{}{ejercicio.3.2}{}}
\newlabel{eje:3-rejilla@cref}{{[ejercicio][2][3]3.2}{[1][21][]21}}
\newlabel{eje:3-cash}{{3.3}{21}{}{ejercicio.3.3}{}}
\newlabel{eje:3-cash@cref}{{[ejercicio][3][3]3.3}{[1][21][]21}}
\newlabel{eje:3-sh}{{3.4}{21}{}{ejercicio.3.4}{}}
\newlabel{eje:3-sh@cref}{{[ejercicio][4][3]3.4}{[1][21][]21}}
\newlabel{eje:3-celdas}{{3.5}{21}{}{ejercicio.3.5}{}}
\newlabel{eje:3-celdas@cref}{{[ejercicio][5][3]3.5}{[1][21][]21}}
\newlabel{eje:3-darts}{{3.6}{21}{}{ejercicio.3.6}{}}
\newlabel{eje:3-darts@cref}{{[ejercicio][6][3]3.6}{[1][21][]21}}
\newlabel{eje:3-coste}{{3.7}{21}{}{ejercicio.3.7}{}}
\newlabel{eje:3-coste@cref}{{[ejercicio][7][3]3.7}{[1][21][]21}}
\newlabel{eje:3-cruce}{{3.8}{22}{}{ejercicio.3.8}{}}
\newlabel{eje:3-cruce@cref}{{[ejercicio][8][3]3.8}{[1][22][]22}}
\newlabel{fig:produccion}{{4.1}{23}{Arquitectura de un optimizador en producción. Las cajas sombreadas son almacenes versionados; las flechas discontinuas, los lazos de reentrenamiento y de consulta de características. GMAO: gestión del mantenimiento asistido por ordenador}{figure.4.1}{}}
\newlabel{fig:produccion@cref}{{[figure][1][4]4.1}{[1][23][]23}}
\newlabel{ej:psi}{{4.7}{25}{Cambio de sensor}{ejemplo.7}{}}
\newlabel{ej:psi@cref}{{[ejemplo][7][]4.7}{[1][25][]25}}
\newlabel{ej:caso-mant}{{4.14}{27}{Optimizador de mantenimiento de una línea de envasado}{ejemplo.14}{}}
\newlabel{ej:caso-mant@cref}{{[ejemplo][14][]4.14}{[1][27][]27}}
\newlabel{obs:errores4}{{4.15}{29}{Errores frecuentes en la práctica}{observacion.15}{}}
\newlabel{obs:errores4@cref}{{[observacion][15][]4.15}{[1][29][]29}}
\newlabel{eje:4-deriva}{{4.1}{30}{}{ejercicio.4.1}{}}
\newlabel{eje:4-deriva@cref}{{[ejercicio][1][4]4.1}{[1][30][]30}}
\newlabel{eje:4-deuda}{{4.2}{30}{}{ejercicio.4.2}{}}
\newlabel{eje:4-deuda@cref}{{[ejercicio][2][4]4.2}{[1][30][]30}}
\newlabel{eje:4-uniforme}{{4.3}{30}{}{ejercicio.4.3}{}}
\newlabel{eje:4-uniforme@cref}{{[ejercicio][3][4]4.3}{[1][30][]30}}
\newlabel{eje:4-psi}{{4.4}{30}{}{ejercicio.4.4}{}}
\newlabel{eje:4-psi@cref}{{[ejercicio][4][4]4.4}{[1][30][]30}}
\newlabel{eje:4-edad}{{4.5}{31}{}{ejercicio.4.5}{}}
\newlabel{eje:4-edad@cref}{{[ejercicio][5][4]4.5}{[1][31][]31}}
\newlabel{eje:4-ips}{{4.6}{31}{}{ejercicio.4.6}{}}
\newlabel{eje:4-ips@cref}{{[ejercicio][6][4]4.6}{[1][31][]31}}
\newlabel{eje:4-reentreno}{{4.7}{31}{}{ejercicio.4.7}{}}
\newlabel{eje:4-reentreno@cref}{{[ejercicio][7][4]4.7}{[1][31][]31}}
\newlabel{eje:4-lineabase}{{4.8}{31}{}{ejercicio.4.8}{}}
\newlabel{eje:4-lineabase@cref}{{[ejercicio][8][4]4.8}{[1][31][]31}}
\newlabel{eje:4-monitor}{{4.9}{31}{}{ejercicio.4.9}{}}
\newlabel{eje:4-monitor@cref}{{[ejercicio][9][4]4.9}{[1][31][]31}}
\newlabel{ej:caso-valor}{{5.6}{33}{Valoración del optimizador del \cref {ej:caso-mant}}{ejemplo.6}{}}
\newlabel{ej:caso-valor@cref}{{[ejemplo][6][]5.6}{[1][33][]33}}
\newlabel{obs:errores5}{{5.9}{36}{Errores frecuentes en la práctica}{observacion.9}{}}
\newlabel{obs:errores5@cref}{{[observacion][9][]5.9}{[1][36][]36}}
\newlabel{eje:5-alternativa}{{5.1}{37}{}{ejercicio.5.1}{}}
\newlabel{eje:5-alternativa@cref}{{[ejercicio][1][5]5.1}{[1][37][]37}}
\newlabel{eje:5-pareado}{{5.2}{37}{}{ejercicio.5.2}{}}
\newlabel{eje:5-pareado@cref}{{[ejercicio][2][5]5.2}{[1][37][]37}}
\newlabel{eje:5-bootstrap}{{5.3}{37}{}{ejercicio.5.3}{}}
\newlabel{eje:5-bootstrap@cref}{{[ejercicio][3][5]5.3}{[1][37][]37}}
\newlabel{eje:5-van}{{5.4}{37}{}{ejercicio.5.4}{}}
\newlabel{eje:5-van@cref}{{[ejercicio][4][5]5.4}{[1][37][]37}}
\newlabel{eje:5-tamano}{{5.5}{37}{}{ejercicio.5.5}{}}
\newlabel{eje:5-tamano@cref}{{[ejercicio][5][5]5.5}{[1][37][]37}}
\newlabel{eje:5-cvar}{{5.6}{38}{}{ejercicio.5.6}{}}
\newlabel{eje:5-cvar@cref}{{[ejercicio][6][5]5.6}{[1][38][]38}}
\newlabel{eje:5-amfe}{{5.7}{38}{}{ejercicio.5.7}{}}
\newlabel{eje:5-amfe@cref}{{[ejercicio][7][5]5.7}{[1][38][]38}}
\newlabel{eje:5-reglamento}{{5.8}{38}{}{ejercicio.5.8}{}}
\newlabel{eje:5-reglamento@cref}{{[ejercicio][8][5]5.8}{[1][38][]38}}
\newlabel{ej:mini-fabrica}{{6.2}{40}{Cambios de formato}{ejemplo.2}{}}
\newlabel{ej:mini-fabrica@cref}{{[ejemplo][2][]6.2}{[1][40][]40}}
\newlabel{ej:mini-3d}{{6.3}{41}{Niveles de detalle con presupuesto}{ejemplo.3}{}}
\newlabel{ej:mini-3d@cref}{{[ejemplo][3][]6.3}{[1][41][]41}}
\newlabel{ej:mini-cpd}{{6.4}{42}{Heurística de colocación y refrigeración}{ejemplo.4}{}}
\newlabel{ej:mini-cpd@cref}{{[ejemplo][4][]6.4}{[1][42][]42}}
\newlabel{obs:errores6}{{6.5}{42}{Errores frecuentes en la práctica}{observacion.5}{}}
\newlabel{obs:errores6@cref}{{[observacion][5][]6.5}{[1][42][]42}}
\newlabel{eje:6-orden}{{6.1}{43}{}{ejercicio.6.1}{}}
\newlabel{eje:6-orden@cref}{{[ejercicio][1][6]6.1}{[1][43][]43}}
\newlabel{eje:6-margenes}{{6.2}{43}{}{ejercicio.6.2}{}}
\newlabel{eje:6-margenes@cref}{{[ejercicio][2][6]6.2}{[1][43][]43}}
\newlabel{eje:6-alternativas}{{6.3}{43}{}{ejercicio.6.3}{}}
\newlabel{eje:6-alternativas@cref}{{[ejercicio][3][6]6.3}{[1][43][]43}}
\newlabel{eje:6-techo}{{6.4}{43}{}{ejercicio.6.4}{}}
\newlabel{eje:6-techo@cref}{{[ejercicio][4][6]6.4}{[1][43][]43}}
\newlabel{eje:6-formatos}{{6.5}{44}{}{ejercicio.6.5}{}}
\newlabel{eje:6-formatos@cref}{{[ejercicio][5][6]6.5}{[1][44][]44}}
\newlabel{eje:6-lod}{{6.6}{44}{}{ejercicio.6.6}{}}
\newlabel{eje:6-lod@cref}{{[ejercicio][6][6]6.6}{[1][44][]44}}
\newlabel{eje:6-simpson}{{6.7}{44}{}{ejercicio.6.7}{}}
\newlabel{eje:6-simpson@cref}{{[ejercicio][7][6]6.7}{[1][44][]44}}
\newlabel{eje:6-diseno}{{6.8}{44}{}{ejercicio.6.8}{}}
\newlabel{eje:6-diseno@cref}{{[ejercicio][8][6]6.8}{[1][44][]44}}
\newlabel{ap:soluciones}{{A}{45}{Soluciones de los ejercicios}{appendix.A}{}}
\newlabel{ap:soluciones@cref}{{[appendix][1][2147483647]A}{[1][45][]45}}
\newlabel{sol:eje:1-componentes}{{A.1}{45}{Tema 1}{section*.3}{}}
\newlabel{sol:eje:1-componentes@cref}{{[subappendix][1][2147483647,1]A.1}{[1][45][]45}}
\newlabel{sol:eje:1-softmax}{{A.1}{45}{Tema 1}{section*.4}{}}
\newlabel{sol:eje:1-softmax@cref}{{[subappendix][1][2147483647,1]A.1}{[1][45][]45}}
\newlabel{sol:eje:1-ffbf}{{A.1}{46}{Tema 1}{section*.5}{}}
\newlabel{sol:eje:1-ffbf@cref}{{[subappendix][1][2147483647,1]A.1}{[1][46][]46}}
\newlabel{sol:eje:1-nextfit}{{A.1}{46}{Tema 1}{section*.6}{}}
\newlabel{sol:eje:1-nextfit@cref}{{[subappendix][1][2147483647,1]A.1}{[1][46][]46}}
\newlabel{sol:eje:1-presupuesto}{{A.1}{47}{Tema 1}{section*.7}{}}
\newlabel{sol:eje:1-presupuesto@cref}{{[subappendix][1][2147483647,1]A.1}{[1][47][]47}}
\newlabel{sol:eje:1-gramatica}{{A.1}{47}{Tema 1}{section*.8}{}}
\newlabel{sol:eje:1-gramatica@cref}{{[subappendix][1][2147483647,1]A.1}{[1][47][]47}}
\newlabel{sol:eje:1-interpretar}{{A.1}{47}{Tema 1}{section*.9}{}}
\newlabel{sol:eje:1-interpretar@cref}{{[subappendix][1][2147483647,1]A.1}{[1][47][]47}}
\newlabel{sol:eje:1-diseno}{{A.1}{48}{Tema 1}{section*.10}{}}
\newlabel{sol:eje:1-diseno@cref}{{[subappendix][1][2147483647,1]A.1}{[1][48][]48}}
\newlabel{sol:eje:2-meta}{{A.2}{49}{Tema 2}{section*.11}{}}
\newlabel{sol:eje:2-meta@cref}{{[subappendix][2][2147483647,1]A.2}{[1][48][]49}}
\newlabel{sol:eje:2-metrica}{{A.2}{49}{Tema 2}{section*.12}{}}
\newlabel{sol:eje:2-metrica@cref}{{[subappendix][2][2147483647,1]A.2}{[1][49][]49}}
\newlabel{sol:eje:2-maldicion}{{A.2}{49}{Tema 2}{section*.13}{}}
\newlabel{sol:eje:2-maldicion@cref}{{[subappendix][2][2147483647,1]A.2}{[1][49][]49}}
\newlabel{sol:eje:2-voto}{{A.2}{50}{Tema 2}{section*.14}{}}
\newlabel{sol:eje:2-voto@cref}{{[subappendix][2][2147483647,1]A.2}{[1][50][]50}}
\newlabel{sol:eje:2-amdahl}{{A.2}{50}{Tema 2}{section*.15}{}}
\newlabel{sol:eje:2-amdahl@cref}{{[subappendix][2][2147483647,1]A.2}{[1][50][]50}}
\newlabel{sol:eje:2-medida}{{A.2}{50}{Tema 2}{section*.16}{}}
\newlabel{sol:eje:2-medida@cref}{{[subappendix][2][2147483647,1]A.2}{[1][50][]50}}
\newlabel{sol:eje:2-cota}{{A.2}{51}{Tema 2}{section*.17}{}}
\newlabel{sol:eje:2-cota@cref}{{[subappendix][2][2147483647,1]A.2}{[1][51][]51}}
\newlabel{sol:eje:2-diseno}{{A.2}{51}{Tema 2}{section*.18}{}}
\newlabel{sol:eje:2-diseno@cref}{{[subappendix][2][2147483647,1]A.2}{[1][51][]51}}
\newlabel{sol:eje:3-dimensiones}{{A.3}{52}{Tema 3}{section*.19}{}}
\newlabel{sol:eje:3-dimensiones@cref}{{[subappendix][3][2147483647,1]A.3}{[1][52][]52}}
\newlabel{sol:eje:3-rejilla}{{A.3}{52}{Tema 3}{section*.20}{}}
\newlabel{sol:eje:3-rejilla@cref}{{[subappendix][3][2147483647,1]A.3}{[1][52][]52}}
\newlabel{sol:eje:3-cash}{{A.3}{52}{Tema 3}{section*.21}{}}
\newlabel{sol:eje:3-cash@cref}{{[subappendix][3][2147483647,1]A.3}{[1][52][]52}}
\newlabel{sol:eje:3-sh}{{A.3}{53}{Tema 3}{section*.22}{}}
\newlabel{sol:eje:3-sh@cref}{{[subappendix][3][2147483647,1]A.3}{[1][53][]53}}
\newlabel{sol:eje:3-celdas}{{A.3}{53}{Tema 3}{section*.23}{}}
\newlabel{sol:eje:3-celdas@cref}{{[subappendix][3][2147483647,1]A.3}{[1][53][]53}}
\newlabel{sol:eje:3-darts}{{A.3}{53}{Tema 3}{section*.24}{}}
\newlabel{sol:eje:3-darts@cref}{{[subappendix][3][2147483647,1]A.3}{[1][53][]53}}
\newlabel{sol:eje:3-coste}{{A.3}{54}{Tema 3}{section*.25}{}}
\newlabel{sol:eje:3-coste@cref}{{[subappendix][3][2147483647,1]A.3}{[1][54][]54}}
\newlabel{sol:eje:3-cruce}{{A.3}{54}{Tema 3}{section*.26}{}}
\newlabel{sol:eje:3-cruce@cref}{{[subappendix][3][2147483647,1]A.3}{[1][54][]54}}
\newlabel{sol:eje:4-deriva}{{A.4}{55}{Tema 4}{section*.27}{}}
\newlabel{sol:eje:4-deriva@cref}{{[subappendix][4][2147483647,1]A.4}{[1][54][]55}}
\newlabel{sol:eje:4-deuda}{{A.4}{55}{Tema 4}{section*.28}{}}
\newlabel{sol:eje:4-deuda@cref}{{[subappendix][4][2147483647,1]A.4}{[1][55][]55}}
\newlabel{sol:eje:4-uniforme}{{A.4}{55}{Tema 4}{section*.29}{}}
\newlabel{sol:eje:4-uniforme@cref}{{[subappendix][4][2147483647,1]A.4}{[1][55][]55}}
\newlabel{sol:eje:4-psi}{{A.4}{56}{Tema 4}{section*.30}{}}
\newlabel{sol:eje:4-psi@cref}{{[subappendix][4][2147483647,1]A.4}{[1][56][]56}}
\newlabel{sol:eje:4-edad}{{A.4}{56}{Tema 4}{section*.31}{}}
\newlabel{sol:eje:4-edad@cref}{{[subappendix][4][2147483647,1]A.4}{[1][56][]56}}
\newlabel{sol:eje:4-ips}{{A.4}{56}{Tema 4}{section*.32}{}}
\newlabel{sol:eje:4-ips@cref}{{[subappendix][4][2147483647,1]A.4}{[1][56][]56}}
\newlabel{sol:eje:4-reentreno}{{A.4}{57}{Tema 4}{section*.33}{}}
\newlabel{sol:eje:4-reentreno@cref}{{[subappendix][4][2147483647,1]A.4}{[1][57][]57}}
\newlabel{sol:eje:4-lineabase}{{A.4}{57}{Tema 4}{section*.34}{}}
\newlabel{sol:eje:4-lineabase@cref}{{[subappendix][4][2147483647,1]A.4}{[1][57][]57}}
\newlabel{sol:eje:4-monitor}{{A.4}{57}{Tema 4}{section*.35}{}}
\newlabel{sol:eje:4-monitor@cref}{{[subappendix][4][2147483647,1]A.4}{[1][57][]57}}
\newlabel{sol:eje:5-alternativa}{{A.5}{58}{Tema 5}{section*.36}{}}
\newlabel{sol:eje:5-alternativa@cref}{{[subappendix][5][2147483647,1]A.5}{[1][58][]58}}
\newlabel{sol:eje:5-pareado}{{A.5}{58}{Tema 5}{section*.37}{}}
\newlabel{sol:eje:5-pareado@cref}{{[subappendix][5][2147483647,1]A.5}{[1][58][]58}}
\newlabel{sol:eje:5-bootstrap}{{A.5}{59}{Tema 5}{section*.38}{}}
\newlabel{sol:eje:5-bootstrap@cref}{{[subappendix][5][2147483647,1]A.5}{[1][59][]59}}
\newlabel{sol:eje:5-van}{{A.5}{59}{Tema 5}{section*.39}{}}
\newlabel{sol:eje:5-van@cref}{{[subappendix][5][2147483647,1]A.5}{[1][59][]59}}
\newlabel{sol:eje:5-tamano}{{A.5}{59}{Tema 5}{section*.40}{}}
\newlabel{sol:eje:5-tamano@cref}{{[subappendix][5][2147483647,1]A.5}{[1][59][]59}}
\newlabel{sol:eje:5-cvar}{{A.5}{60}{Tema 5}{section*.41}{}}
\newlabel{sol:eje:5-cvar@cref}{{[subappendix][5][2147483647,1]A.5}{[1][60][]60}}
\newlabel{sol:eje:5-amfe}{{A.5}{60}{Tema 5}{section*.42}{}}
\newlabel{sol:eje:5-amfe@cref}{{[subappendix][5][2147483647,1]A.5}{[1][60][]60}}
\newlabel{sol:eje:5-reglamento}{{A.5}{60}{Tema 5}{section*.43}{}}
\newlabel{sol:eje:5-reglamento@cref}{{[subappendix][5][2147483647,1]A.5}{[1][60][]60}}
\newlabel{sol:eje:6-orden}{{A.6}{61}{Tema 6}{section*.44}{}}
\newlabel{sol:eje:6-orden@cref}{{[subappendix][6][2147483647,1]A.6}{[1][61][]61}}
\newlabel{sol:eje:6-margenes}{{A.6}{61}{Tema 6}{section*.45}{}}
\newlabel{sol:eje:6-margenes@cref}{{[subappendix][6][2147483647,1]A.6}{[1][61][]61}}
\newlabel{sol:eje:6-alternativas}{{A.6}{62}{Tema 6}{section*.46}{}}
\newlabel{sol:eje:6-alternativas@cref}{{[subappendix][6][2147483647,1]A.6}{[1][62][]62}}
\newlabel{sol:eje:6-techo}{{A.6}{62}{Tema 6}{section*.47}{}}
\newlabel{sol:eje:6-techo@cref}{{[subappendix][6][2147483647,1]A.6}{[1][62][]62}}
\newlabel{sol:eje:6-formatos}{{A.6}{62}{Tema 6}{section*.48}{}}
\newlabel{sol:eje:6-formatos@cref}{{[subappendix][6][2147483647,1]A.6}{[1][62][]62}}
\newlabel{sol:eje:6-lod}{{A.6}{63}{Tema 6}{section*.49}{}}
\newlabel{sol:eje:6-lod@cref}{{[subappendix][6][2147483647,1]A.6}{[1][63][]63}}
\newlabel{sol:eje:6-simpson}{{A.6}{63}{Tema 6}{section*.50}{}}
\newlabel{sol:eje:6-simpson@cref}{{[subappendix][6][2147483647,1]A.6}{[1][63][]63}}
\newlabel{sol:eje:6-diseno}{{A.6}{64}{Tema 6}{section*.51}{}}
\newlabel{sol:eje:6-diseno@cref}{{[subappendix][6][2147483647,1]A.6}{[1][64][]64}}
\makeatother

\begin{document}
\frontmatter
\begin{titlepage}
\centering
\vspace*{2.5cm}
{\large\scshape Anatomía de un LLM}\par
\vspace{0.4cm}
{\normalsize Fundamentos matemáticos de los modelos de lenguaje}\par
\vspace{2.2cm}
\rule{\textwidth}{0.8pt}\par\vspace{0.5cm}
{\small\scshape Fase VI \textperiodcentered\ Diseño y construcción de optimizadores con IA}\par\vspace{0.3cm}
{\Huge\bfseries Módulo 28\par}\vspace{0.3cm}
{\LARGE LLM como optimizadores y del prototipo a la producción\par}
\vspace{0.2cm}
{\large Búsqueda evolutiva de programas, optimización de sistemas con LLM, AutoML, MLOps, valor y riesgo, y guía de diseño\par}
\vspace{0.5cm}\rule{\textwidth}{0.4pt}\par
\vspace{1.6cm}
{\large Resumen del módulo}\par
\vfill
{\small Edición de octubre de 2026}\par
\end{titlepage}

\chapter*{Sobre este resumen}
\addcontentsline{toc}{chapter}{Sobre este resumen}
\markboth{Sobre este resumen}{}

Este documento reúne lo indispensable de los apuntes del Módulo 28: el texto expositivo, todas las definiciones y todos los resultados (teoremas, proposiciones, lemas y corolarios), las tablas y las observaciones de contenido. Se han suprimido las demostraciones, los ejemplos, las lecturas, las figuras, los ejercicios y sus soluciones. La numeración de temas, secciones, enunciados y ecuaciones es la misma que en los apuntes, de modo que cualquier referencia puede buscarse allí; las que remiten a elementos suprimidos (un ejemplo, una figura, una ecuación de una demostración) conservan su número en los apuntes.

\section*{Notación}
\begin{center}\small
\begin{tabularx}{\textwidth}{@{}l>{\raggedright\arraybackslash}X@{}}
\toprule
Símbolo & Significado \\
\midrule
$\Prog$, $\pi\in\Prog$ & espacio de programas (o de instrucciones, configuraciones) y un elemento \\
$h(\pi)$, $\hat h(\pi)$ & puntuación verdadera del programa y puntuación medida por el evaluador \\
$q$ & probabilidad de que una propuesta mejore al mejor programa actual \\
$C$, $s_j$, $r$, $g=r-s$ & capacidad del contenedor, tamaño del objeto $j$, capacidad restante y hueco tras colocar \\
$L_1=\lceil\sum_j s_j/C\rceil$ & cota inferior del número óptimo de contenedores \\
$K$, $n$ & número de candidatos comparados y tamaño del conjunto de validación \\
$\eta$, $r_{\min}$, $r_{\max}$ & factor de reducción y presupuestos mínimo y máximo de la multifidelidad \\
$R(t)$, $F(t)$, $f(t)$, $\lambda(t)$ & fiabilidad, función de distribución, densidad y tasa de fallo de una vida útil \\
$\beta$, $\theta$ & parámetros de forma y de escala de la distribución de Weibull \\
$c_p$, $c_f$ & coste de una intervención preventiva y de una avería \\
$\PSI$, $\KL$ & índice de estabilidad poblacional y divergencia de Kullback--Leibler \\
$\VAN$, $\rho$ & valor actual neto y tipo de descuento anual \\
$\mathrm{CVaR}_\alpha$ & valor en riesgo condicional al nivel $\alpha$ \\
€, h, d, mm/s, ºC & euros, horas, días, milímetros por segundo (vibración) y grados Celsius \\
\bottomrule
\end{tabularx}
\end{center}


\tableofcontents

\mainmatter
% =====================================================================
\chapter{Búsqueda evolutiva de programas guiada por LLM}
\label{tema:1}
% =====================================================================

En muchos problemas lo que se busca no es un número ni un vector, sino un \emph{procedimiento}: una regla de prioridad para un planificador, una heurística de empaquetado, un núcleo de cálculo más rápido. Este tema estudia cómo un modelo de lenguaje, encerrado en un bucle evolutivo con un evaluador automático, puede buscar en el espacio de programas. Sigue el artículo de FunSearch de Romera-Paredes \emph{et al.}~\cite{funsearch} y su implementación de referencia~\cite{funsearch-code}, el informe de AlphaEvolve de Novikov \emph{et al.}~\cite{alphaevolve} y la nota técnica de Google DeepMind~\cite{ae-blog}; para la representación de programas como expresiones se apoya en el capítulo «Expression Optimization» de Kochenderfer y Wheeler~\cite{kw}, y para el empaquetado, en la §3.3 de Williamson y Shmoys~\cite{ws}. La lectura crítica de Herrmann y Pallez~\cite{hp} sirve de contrapeso.

\section{Optimizar en el espacio de programas}

\begin{definicion}[Búsqueda de programas]\label{def:busqueda-programas}
Sea $\Prog$ el conjunto de programas que respetan una firma dada (por ejemplo, funciones que reciben un objeto y la lista de contenedores abiertos y devuelven una puntuación por contenedor). Sea $\mathcal{D}$ una distribución de instancias del problema y $u(\pi,x)\in\R\cup\{-\infty\}$ la utilidad de ejecutar $\pi$ sobre la instancia $x$, con $u=-\infty$ si el programa falla, excede el tiempo o produce una salida inválida. La \emph{búsqueda de programas} es
\[
\max_{\pi\in\Prog}\ h(\pi),\qquad h(\pi)=\E_{x\sim\mathcal{D}}\bigl[u(\pi,x)\bigr].
\]
Como $\mathcal{D}$ no se conoce, se usa un \emph{evaluador} $\hat h(\pi)=\frac1m\sum_{i=1}^{m}u(\pi,x_i)$ sobre instancias de entrenamiento $x_1,\dots,x_m$, y la calidad final se mide con instancias de prueba independientes.
\end{definicion}

El problema es de caja negra en el sentido de la definición~5.1 del módulo~24: no hay gradiente, el espacio es discreto y enorme, y cada evaluación cuesta una ejecución. La programación genética clásica busca en ese espacio con mutaciones y cruces aleatorios de árboles sintácticos~\cite{kw}; su dificultad es que casi todas las mutaciones aleatorias producen programas inútiles. La idea de FunSearch es sustituir la mutación aleatoria por un modelo de lenguaje entrenado con código, que propone modificaciones \emph{plausibles}.

\begin{observacion}[Por qué buscar programas y no soluciones]\label{obs:por-que-programas}
(i) Un programa \emph{generaliza}: una heurística de empaquetado sirve para instancias nuevas, mientras que un empaquetado concreto solo sirve para la suya. (ii) Un programa corto puede \emph{describir} un objeto enorme: en el problema del \emph{cap set}, FunSearch buscó programas que construyen el conjunto, no el conjunto~\cite{funsearch}. (iii) Un programa se puede \emph{leer}: el experto puede entender, simplificar y auditar la regla encontrada. (iv) Un programa se puede \emph{verificar}: el evaluador lo ejecuta y comprueba la validez de la salida.
\end{observacion}

\begin{definicion}[Esqueleto y función que evoluciona]\label{def:esqueleto}
Un \emph{esqueleto} es un programa fijo que resuelve el problema llamando a una función con firma conocida, la única que se modifica durante la búsqueda. En el empaquetado en línea, el esqueleto recorre los objetos en orden de llegada, calcula los contenedores en los que cabe cada uno y coloca el objeto en el de mayor \texttt{prioridad(objeto, contenedores)}; solo \texttt{prioridad} evoluciona.
\end{definicion}

El esqueleto encierra el conocimiento del problema y garantiza la factibilidad (un objeto nunca se coloca donde no cabe), de modo que cualquier función que devuelva números produce una solución válida. Es el mismo principio de diseño que el filtro de seguridad del módulo~26 (tema~6, §6.5): la parte aprendida propone dentro de un marco que impide lo inaceptable.

\section{FunSearch}

\addtocounter{figure}{1}%

El bucle de la \cref{fig:funsearch} se ejecuta de forma asíncrona y distribuida. Según el artículo y la implementación publicada~\cite{funsearch,funsearch-code}:
\begin{enumerate}[label=\arabic*.]
  \item \emph{Constructor del prompt.} Se elige una isla al azar y, dentro de ella, $k=2$ programas; se ordenan de peor a mejor puntuación, se renombran como \texttt{prioridad\_v0} y \texttt{prioridad\_v1}, y el prompt termina con la cabecera de \texttt{prioridad\_v2}, «versión mejorada de la v1». A esta técnica los autores la llaman \emph{best-shot prompting}.
  \item \emph{Muestreo.} Un LLM preentrenado con código (en el trabajo original, un modelo de la familia PaLM~2; en la actualización de 2024, Gemini~1.5 Flash~\cite{fs-blog}) completa el cuerpo de la nueva función. La configuración por defecto de la implementación usa 15 muestreadores que piden 4 continuaciones por prompt.
  \item \emph{Evaluación.} El programa se ejecuta en un entorno aislado sobre las instancias de entrenamiento; si falla, tarda demasiado o produce algo inválido, se descarta. La configuración por defecto usa 140 evaluadores en paralelo.
  \item \emph{Base de programas.} Los programas válidos se guardan en $10$ islas que evolucionan por separado. Dentro de cada isla se agrupan en \emph{clusters} por su \emph{firma} (el vector de puntuaciones en cada instancia de entrenamiento). Periódicamente se reinicia la mitad de las islas con peor mejor puntuación, sembrando cada una con el mejor programa de una isla superviviente elegida al azar.
\end{enumerate}

\begin{definicion}[Muestreo de clusters con temperatura]\label{def:softmax-clusters}
Sean $s_1,\dots,s_c$ las puntuaciones de los clusters de una isla y $T>0$ una temperatura. Se elige el cluster $i$ con probabilidad
\[
p_i(T)=\frac{\exp(s_i/T)}{\sum_{l=1}^{c}\exp(s_l/T)}.
\]
En la implementación de referencia, $T$ decrece linealmente con el número $n$ de programas registrados en la isla, $T=T_0\bigl(1-(n\bmod P)/P\bigr)$, con $T_0=0{,}1$ y periodo $P=30\,000$; dentro del cluster elegido se prefieren los programas más cortos, con una softmax sobre la longitud normalizada cambiada de signo~\cite{funsearch-code}.
\end{definicion}

\begin{proposicion}[Límites de la temperatura]\label{prop:temperatura}
Sea $M=\max_i s_i$ y $A=\{i: s_i=M\}$. Cuando $T\to 0^{+}$, $p_i(T)\to 1/\abs{A}$ si $i\in A$ y $p_i(T)\to 0$ en otro caso. Cuando $T\to\infty$, $p_i(T)\to 1/c$ para todo $i$.
\end{proposicion}

La temperatura regula, pues, el equilibrio entre \emph{explotar} los mejores programas y \emph{explorar} otros que pueden ser punto de partida de mejoras mayores; es la misma distribución de Boltzmann del recocido simulado (módulo~24, tema~5, §5.2). Las islas y los reinicios cumplen una función parecida a otra escala: mantienen poblaciones que no se contaminan entre sí y, de vez en cuando, abandonan las que se han estancado.

\begin{proposicion}[Elitismo y espera de una mejora]\label{prop:espera}
Sea $b_t$ la mejor puntuación de la base de programas tras $t$ evaluaciones, en un algoritmo que nunca elimina el mejor programa. Entonces:
\begin{enumerate}
  \item $b_{t+1}\ge b_t$ para todo $t$.
  \item Si, mientras no haya mejora, cada propuesta supera a $b_t$ con probabilidad al menos $q>0$, de forma independiente, el número $N$ de evaluaciones hasta la primera mejora cumple $\E[N]\le 1/q$ y $\Prob(N>B)\le(1-q)^{B}\le e^{-qB}$.
  \item Si $I$ islas independientes tienen cada una probabilidad $p$ de alcanzar una puntuación objetivo con su presupuesto, la probabilidad de que la alcance alguna es $1-(1-p)^{I}$.
\end{enumerate}
\end{proposicion}

La proposición explica el papel del LLM con una sola cifra: el presupuesto necesario es del orden de $1/q$. Una mutación aleatoria de un árbol sintáctico tiene una $q$ minúscula; un modelo que ha visto millones de programas concentra la probabilidad en ediciones que tienen sentido (cambiar un umbral, añadir un caso, simplificar una expresión), y eso multiplica $q$.

\section{AlphaEvolve}

AlphaEvolve~\cite{alphaevolve} generaliza FunSearch de función a sistema. El usuario marca con comentarios (\texttt{EVOLVE-BLOCK-START} y \texttt{EVOLVE-BLOCK-END}) los bloques de un fichero que pueden cambiar; el resto actúa como esqueleto. Los LLM devuelven \emph{diferencias} en formato de búsqueda y sustitución en lugar de reescribir el fichero, lo que permite evolucionar programas de cientos de líneas en cualquier lenguaje. Un conjunto de modelos combina uno rápido (Gemini~2.0 Flash, muchas propuestas) y uno más capaz (Gemini~2.0 Pro, menos propuestas y de más calidad). La evaluación es una \emph{cascada}: el programa pasa por pruebas de dificultad creciente y solo avanza si supera las anteriores, de modo que los candidatos malos se descartan con poco coste. La base de programas combina MAP-Elites (una población que guarda el mejor programa en cada celda de un espacio de características, para mantener diversidad) con el modelo de islas. El \cref{tab:fs-ae} resume las diferencias que los propios autores señalan.

\begin{table}[h]
\centering\small
\caption{FunSearch frente a AlphaEvolve, según la tabla~1 de~\cite{alphaevolve}.}\label{tab:fs-ae}
\begin{tabularx}{\textwidth}{@{}>{\raggedright\arraybackslash}p{3.0cm}>{\raggedright\arraybackslash}X>{\raggedright\arraybackslash}X@{}}
\toprule
 & FunSearch & AlphaEvolve \\
\midrule
qué evoluciona & una función & un fichero entero, con varios bloques \\
tamaño & 10--20 líneas & hasta cientos de líneas \\
lenguaje & Python & cualquiera \\
evaluación & $\le 20$~min en una CPU & horas, en paralelo, con aceleradores \\
muestras del LLM & millones & bastan miles \\
\bottomrule
\end{tabularx}
\end{table}

\addtocounter{definicion}{1}%

\section{Por qué funciona y cuándo no}

Tres ingredientes explican el éxito. \emph{Asimetría entre generar y verificar}: comprobar que un empaquetado es válido y contar contenedores es inmediato, aunque encontrar una buena heurística sea difícil; es la misma asimetría que explota el aprendizaje por refuerzo con recompensas verificables (módulo~21, tema~6, §6.2) y el esquema de generar y verificar del módulo~27 (tema~6, §6.3). \emph{Una distribución de propuestas informada}: el LLM eleva la probabilidad $q$ de la \cref{prop:espera}. \emph{Memoria y diversidad}: la base de programas acumula lo aprendido y las islas evitan la convergencia prematura.

Los límites son igual de claros. (i) \emph{El evaluador es la especificación}: lo que no mida no se optimiza, y lo que mida mal se explota, como en la sobreoptimización de la recompensa (módulo~21, tema~6, §6.4). (ii) \emph{Generalización}: la heurística se ajusta a la distribución de entrenamiento; Herrmann y Pallez~\cite{hp} concluyen que las heurísticas de empaquetado de FunSearch no son robustas a cambios en la distribución de los datos y que en gran parte redescubren patrones conocidos. (iii) \emph{Coste}: millones de muestras en FunSearch y horas de evaluación en AlphaEvolve. (iv) \emph{Seguridad}: el código generado es código no confiable y debe ejecutarse aislado, sin red ni acceso a datos sensibles.

\section{Empaquetado en línea}

\begin{definicion}[Empaquetado unidimensional en línea]\label{def:bp}
Se reciben, uno a uno, objetos de tamaños $s_1,s_2,\dots,s_N\in(0,C]$ que deben colocarse en contenedores de capacidad $C$ sin superarla. Cada objeto se coloca al llegar, sin conocer los siguientes, y no se puede mover después. El objetivo es minimizar el número de contenedores usados. Sea $\mathrm{OPT}$ el mínimo que se lograría conociendo todos los objetos de antemano.
\end{definicion}

\begin{definicion}[Reglas de prioridad clásicas]\label{def:ff-bf}
Sea $r$ la capacidad restante de un contenedor abierto en el que cabe el objeto $s$ y $g=r-s$ el hueco que dejaría. \emph{First fit} (FF) coloca el objeto en el primer contenedor abierto en que cabe; \emph{best fit} (BF), en el que deja menor hueco $g$; en ambos casos se abre uno nuevo solo si no cabe en ninguno. En términos del esqueleto, FF usa la prioridad $-(\text{índice del contenedor})$ y BF la prioridad $-g$. Un algoritmo es de tipo \emph{any fit} si nunca abre un contenedor nuevo cuando el objeto cabe en uno abierto.
\end{definicion}

\begin{lema}[Cota $L_1$]\label{lem:l1}
$\mathrm{OPT}\ge L_1=\bigl\lceil\sum_{j}s_j/C\bigr\rceil$.
\end{lema}

\begin{teorema}[Garantía de los algoritmos \emph{any fit}]\label{teo:anyfit}
Todo algoritmo \emph{any fit}, en particular FF y BF, usa a lo sumo $\lceil 2\sum_j s_j/C\rceil\le 2\,\mathrm{OPT}$ contenedores.
\end{teorema}

La constante~2 no es la mejor posible para FF, pero el resultado ilustra el tipo de garantía que una heurística generada \emph{no} trae consigo: lo que se sabe de ella es lo que mide el evaluador. Para el problema sin restricciones de llegada, Williamson y Shmoys demuestran que, salvo que $\mathrm{P}=\mathrm{NP}$, ningún algoritmo de tiempo polinómico garantiza un factor menor que $3/2$~\cite[§3.3]{ws}.

\addtocounter{definicion}{1}%

\section{Caso de estudio: una heurística de empaquetado generada}

\addtocounter{definicion}{1}%

\begin{table}[h]
\centering\small
\caption{Empaquetado en línea de 5 instancias de prueba de 5000 objetos ($\bar L_1=1985{,}6$). Datos del cuaderno.}\label{tab:bp}
\begin{tabular}{@{}lcc@{}}
\toprule
Regla & Bandejas (media) & Exceso sobre $L_1$ \\
\midrule
\emph{first fit} & 2073{,}4 & 4{,}42\,\% \\
\emph{best fit} & 2064{,}8 & 3{,}99\,\% \\
regla evolucionada $\pi^{*}$ & 2011{,}2 & 1{,}29\,\% \\
heurística publicada por FunSearch & 2000{,}0 & 0{,}73\,\% \\
\bottomrule
\end{tabular}
\end{table}

Como referencia externa, en el conjunto Weibull de 5000 objetos que acompaña al artículo, BF usa de media 2067{,}0 bandejas y la heurística de FunSearch 2001{,}4, con $\bar L_1=1987{,}8$: excesos del $3{,}98\,\%$ y del $0{,}68\,\%$~\cite{funsearch-code}. La \emph{prueba de robustez} del cuaderno repite la comparación con otra distribución, Weibull de escala 60 (objetos más grandes): BF deja un exceso del $5{,}19\,\%$, la regla $\pi^{*}$ un $5{,}21\,\%$ y la heurística publicada un $6{,}40\,\%$, peor que FF ($6{,}05\,\%$). Con seis semillas distintas del generador, el exceso en prueba de la regla encontrada varía entre el $1{,}29\,\%$ y el $2{,}33\,\%$, y cuatro de las seis reglas tienen la misma estructura: un término de \emph{best fit} saturado más una penalización de los huecos intermedios. Las dos heurísticas generadas se han especializado en la distribución con la que se evaluaron, tal como advierten Herrmann y Pallez~\cite{hp}.

\begin{table}[h]
\centering\small
\caption{Tecnologías del tema~1.}\label{tab:tec1}
\begin{tabularx}{\textwidth}{@{}>{\raggedright\arraybackslash}p{2.5cm}>{\raggedright\arraybackslash}X>{\raggedright\arraybackslash}p{2.1cm}>{\raggedright\arraybackslash}X@{}}
\toprule
Herramienta & Para qué sirve & Licencia o tipo & Cuándo elegirla \\
\midrule
FunSearch (repositorio)~\cite{funsearch-code} & algoritmo evolutivo, manipulación de código y una versión de un solo hilo, sin LLM ni entorno aislado & abierta (Apache~2.0) & entender el método y adaptarlo a un LLM y una infraestructura propios \\
OpenEvolve~\cite{openevolve} & implementación abierta inspirada en AlphaEvolve: islas, MAP-Elites, evaluación en cascada, conjuntos de LLM & abierta (Apache~2.0) & evolucionar ficheros completos con cualquier API compatible \\
AlphaEvolve~\cite{alphaevolve} & agente de Google DeepMind & servicio, no publicado como código & referencia de diseño; resultados publicados \\
gplearn, DEAP & programación genética clásica sin LLM & abiertas (BSD-3; LGPL) & línea base barata para expresiones pequeñas \\
Contenedores aislados & ejecutar código no confiable sin red y con límites de tiempo y memoria & abierta o del proveedor & siempre que se ejecute código generado \\
\bottomrule
\end{tabularx}
\end{table}

\addtocounter{definicion}{1}%


% =====================================================================
\chapter{Optimización de instrucciones, de sistemas con LLM y de código}
\label{tema:2}
% =====================================================================

Un sistema construido con modelos de lenguaje tiene muchos parámetros que no son números: las instrucciones de cada llamada, los ejemplos que se incluyen, el orden de los pasos, el número de muestras que se votan. Este tema los trata como variables de decisión de un problema de optimización de caja negra con un objetivo medido sobre datos. Sigue a Yang \emph{et al.} para OPRO~\cite{opro}, a Khattab \emph{et al.} para DSPy~\cite{dspy} y a Opsahl-Ong \emph{et al.} para MIPRO~\cite{mipro}; la autoconsistencia, a Wang \emph{et al.}~\cite{wang-sc} y al módulo~22 (tema~1, §1.4); el sesgo del ganador, a Smith y Winkler~\cite{smith}; la optimización de código y la medida de tiempos, a la clase~10 «Measurement and Timing» del curso 6.172 del MIT~\cite{mit6172}.

\section{Instrucciones como variables de decisión}

\begin{definicion}[Optimización de instrucciones]\label{def:prompt-opt}
Sea $\mathcal{I}$ un conjunto de instrucciones (textos), $M$ un modelo de lenguaje fijo, $D_{\mathrm{val}}=\{(x_i,y_i)\}_{i=1}^{n}$ un conjunto de validación y $m(\hat y,y)\in[0,1]$ una métrica (por ejemplo, $\Ind[\hat y=y]$). El problema es
\[
\max_{p\in\mathcal{I}}\ J(p),\qquad J(p)=\E_{(x,y)}\bigl[m(M(p,x),y)\bigr],\qquad \hat J(p)=\frac1n\sum_{i=1}^{n}m\bigl(M(p,x_i),y_i\bigr).
\]
Es un problema de caja negra (módulo~24, definición~5.1), discreto, ruidoso (el estimador $\hat J$ tiene varianza $J(1-J)/n$ si $m$ es binaria) y caro (cada evaluación son $n$ llamadas al modelo).
\end{definicion}

\begin{definicion}[OPRO]\label{def:opro}
\emph{Optimization by PROmpting}~\cite{opro} usa un segundo LLM como optimizador. En cada paso se le da un \emph{meta-prompt} con: una descripción de la tarea y unos pocos ejemplos, y la \emph{trayectoria}: las mejores instrucciones encontradas hasta el momento con su puntuación, ordenadas de menor a mayor. El LLM optimizador genera nuevas instrucciones; se evalúan con el LLM que resuelve la tarea, se añaden a la trayectoria y se repite. En los experimentos del artículo, la trayectoria conserva las 20 mejores, se generan 8 instrucciones por paso con temperatura~1 y la puntuación es la exactitud sobre un subconjunto fijo del conjunto de entrenamiento.
\end{definicion}

OPRO es una búsqueda evolutiva como la del \cref{tema:1}, con dos diferencias: el «programa» es un texto y la memoria es la propia trayectoria, que el LLM lee para inferir qué rasgos se asocian a puntuaciones altas. Los autores informan de que las mejores instrucciones encontradas superan a las diseñadas por personas hasta en 8 puntos en GSM8K y hasta en 50 en tareas de Big-Bench Hard~\cite{opro}; también advierten de que, en el problema del viajante, el LLM se mantiene cerca de las heurísticas clásicas con 50 ciudades, pero en instancias mayores la inserción más lejana supera a todos los modelos probados.

\begin{definicion}[Programa de modelos de lenguaje y su compilación]\label{def:dspy}
En DSPy~\cite{dspy}, un \emph{programa de LM} es un grafo de \emph{módulos}; cada módulo tiene una \emph{firma} declarativa (entradas y salidas con nombre, por ejemplo \texttt{pregunta, contexto -> respuesta}) y unos parámetros: su instrucción y sus demostraciones. Un \emph{optimizador} (antes llamado \emph{teleprompter}) recibe el programa, un conjunto de entrenamiento y una métrica, y elige los parámetros de todos los módulos para maximizar la métrica del programa completo. Por ejemplo, el arranque de demostraciones ejecuta el programa sobre los ejemplos de entrenamiento y conserva como demostraciones las trazas que superan la métrica.
\end{definicion}

Khattab \emph{et al.} informan de que los programas compilados superan a los mismos programas con pocos ejemplos estándar en más de un 25\,\% con GPT-3.5 y un 65\,\% con llama2-13b-chat~\cite{dspy}. MIPRO~\cite{mipro} optimiza conjuntamente instrucciones y demostraciones de cada módulo sin etiquetas intermedias: propone instrucciones teniendo en cuenta los datos y el programa, y usa evaluaciones en minilotes para ajustar un modelo sustituto del objetivo, como la optimización bayesiana del módulo~25 (tema~4); obtiene mejoras de hasta 13 puntos de exactitud con Llama-3-8B.

\section{El sesgo del ganador}

Todas estas técnicas eligen el máximo de muchas estimaciones ruidosas. El resultado es predecible.

\begin{proposicion}[Maldición del optimizador]\label{prop:maldicion}
Sean $K$ candidatos con valores verdaderos $J_1,\dots,J_K$ y estimaciones $\hat J_k=J_k+\varepsilon_k$ con $\E[\varepsilon_k]=0$, y sea $\hat k=\argmax_k\hat J_k$. Entonces:
\begin{enumerate}
  \item $\E[\hat J_{\hat k}]\ \ge\ \max_k J_k\ \ge\ \E[J_{\hat k}]$: la puntuación de validación del elegido sobrestima en media su valor verdadero.
  \item Si cada $\varepsilon_k\sim\mathcal{N}(0,\sigma^{2})$ (no necesariamente independientes), $\E[\max_k\varepsilon_k]\le\sigma\sqrt{2\ln K}$. En particular, si todos los $J_k$ son iguales, el sesgo $\E[\hat J_{\hat k}-J_{\hat k}]$ está entre $0$ y $\sigma\sqrt{2\ln K}$.
\end{enumerate}
\end{proposicion}

\addtocounter{definicion}{1}%

\begin{corolario}[Protocolo de evaluación]\label{cor:protocolo}
Para estimar sin sesgo el valor de un sistema optimizado se necesita un conjunto de \emph{prueba} que no haya intervenido en ninguna elección (ni de instrucciones, ni de ejemplos, ni de hiperparámetros), igual que en la separación en entrenamiento, validación y prueba del módulo~8 (tema~1, §1.4). Su error típico, $\sqrt{J(1-J)/n_{\mathrm{test}}}$, fija la resolución con la que se pueden comparar sistemas.
\end{corolario}

\section{Caso de estudio: clasificación de partes de avería}

\addtocounter{definicion}{1}%

\section{La autoconsistencia como problema de optimización}

La autoconsistencia (módulo~22, definición~1.14) devuelve la respuesta más votada entre $m$ muestras. Más muestras dan más exactitud y más coste; el número de muestras es una variable de decisión.

\begin{proposicion}[Incremento de la exactitud del voto]\label{prop:incremento}
Sea $P_m$ la probabilidad de que la mayoría de $m$ (impar) respuestas independientes sea correcta cuando cada una lo es con probabilidad $p>1/2$ (módulo~22, proposición~1.16). Entonces
\[
\Delta_m:=P_{m+2}-P_m=\binom{m}{\frac{m-1}{2}}\bigl(p(1-p)\bigr)^{\frac{m+1}{2}}(2p-1)>0,\qquad \frac{\Delta_{m+2}}{\Delta_m}=\frac{4(m+2)}{m+3}\,p(1-p)<1 .
\]
\end{proposicion}

\begin{corolario}[Número óptimo de muestras]\label{cor:m-optimo}
Si una respuesta correcta vale $V$ y cada muestra cuesta $c$, el beneficio neto $V P_m-cm$ (con $m$ impar) se maximiza en el primer $m$ tal que $V\Delta_m<2c$.
\end{corolario}

Con $p=0{,}65$, $V=0{,}50$~€ y $c=0{,}01$~€: $\Delta_1=0{,}0683$, $\Delta_3=0{,}0466$ y $\Delta_5=0{,}0353$; $V\Delta_3=0{,}0233\ge0{,}02$ pero $V\Delta_5=0{,}0177<0{,}02$, de modo que $m^{*}=5$, con exactitud $P_5=0{,}765$ y beneficio neto $0{,}332$~€ por parte. La hipótesis de independencia es optimista (módulo~22, ejemplo~1.17): conviene estimar $P_m$ con datos y aplicar el mismo criterio.

\section{Optimización de código con agentes}

Un agente de programación (módulo~22, tema~6) puede usarse como optimizador de un programa existente. El planteamiento correcto es el de un problema con restricciones:
\[
\min_{\pi}\ t(\pi)\quad\text{s.a.}\quad \pi \text{ supera todas las pruebas},\quad \text{memoria}(\pi)\le \bar M,\quad \text{coste del agente}\le B,
\]
donde $t$ es el tiempo de ejecución en una carga representativa. Las pruebas son una \emph{restricción dura}, no un término del objetivo: un programa más rápido que falla una prueba no es una solución. Es el mismo principio del evaluador del \cref{tema:1}.

\begin{proposicion}[Ley de Amdahl]\label{prop:amdahl}
Si una fracción $f$ del tiempo de ejecución se acelera en un factor $a>1$ y el resto no cambia, la aceleración global es
\[
A=\frac{1}{(1-f)+f/a}<\frac{1}{1-f}.
\]
\end{proposicion}

La ley dice dónde empezar: por lo que el \emph{perfilador} muestra que ocupa más tiempo. Un agente que no perfila tiende a optimizar lo que parece lento en el código, no lo que lo es.

\addtocounter{definicion}{1}%

\begin{table}[h]
\centering\small
\caption{Tecnologías del tema~2.}\label{tab:tec2}
\begin{tabularx}{\textwidth}{@{}>{\raggedright\arraybackslash}p{2.5cm}>{\raggedright\arraybackslash}X>{\raggedright\arraybackslash}p{2.1cm}>{\raggedright\arraybackslash}X@{}}
\toprule
Herramienta & Para qué sirve & Licencia o tipo & Cuándo elegirla \\
\midrule
DSPy~\cite{dspy} & programas de LM declarativos y optimizadores de instrucciones y demostraciones (incluido MIPRO) & abierta (MIT) & sistemas de varias llamadas con una métrica clara \\
OPRO~\cite{opro} & método: un LLM como optimizador a partir de la trayectoria & artículo con código de los autores & instrucciones de una sola llamada; problemas pequeños \\
promptfoo & conjuntos de pruebas y comparación de instrucciones y modelos & abierta (MIT) & evaluación de regresión de instrucciones \\
pytest, pytest-benchmark & pruebas como restricción y medidas de tiempo repetidas & abiertas (MIT; BSD-2) & optimización de código en Python \\
perf, py-spy, Scalene & perfilado de CPU y memoria & abiertas (GPL-2.0; MIT; Apache~2.0) & antes de pedir al agente que optimice nada \\
asv & seguimiento del rendimiento a lo largo del historial del código & abierta (BSD-3) & evitar regresiones de rendimiento \\
\bottomrule
\end{tabularx}
\end{table}

\addtocounter{definicion}{1}%


% =====================================================================
\chapter{AutoML y búsqueda de arquitecturas}
\label{tema:3}
% =====================================================================

Elegir un modelo, sus hiperparámetros y su preprocesado es también un problema de optimización sobre programas: cada configuración es una canalización ejecutable cuyo valor se mide con validación. Este tema formaliza ese problema siguiendo el libro editado por Hutter, Kotthoff y Vanschoren~\cite{automl} (cap.~1 «Hyperparameter Optimization» y cap.~3 «Neural Architecture Search», de Elsken, Metzen y Hutter~\cite{elsken}) y los artículos originales: Auto-WEKA~\cite{autoweka}, búsqueda aleatoria~\cite{bergstra}, búsqueda de arquitecturas con refuerzo~\cite{zoph}, con evolución~\cite{real}, con pesos compartidos~\cite{enas} y diferenciable~\cite{darts}, y la línea base de Li y Talwalkar~\cite{lt}. La multifidelidad sigue a Jamieson y Talwalkar~\cite{sh}, a Li \emph{et al.}~\cite{hyperband} y a la §19.4 de Zhang \emph{et al.}~\cite{d2l}. Complementa el tema~6 del módulo~25, que estudió el ajuste de hiperparámetros con optimización bayesiana y Hyperband; aquí el énfasis está en el espacio de búsqueda y en el coste.

\section{AutoML como optimización sobre programas}

\begin{definicion}[Selección conjunta de algoritmo e hiperparámetros]\label{def:cash}
Se consideran $R$ algoritmos de aprendizaje, $A^{(1)},\dots,A^{(R)}$, cada uno con un espacio de hiperparámetros $\Lambda^{(j)}$, y sea $D$ un conjunto de datos dividido en $k$ pares $(D^{(i)}_{\mathrm{tr}},D^{(i)}_{\mathrm{va}})$, cada uno formado por una parte de entrenamiento y otra de validación. El problema CASH (\emph{combined algorithm selection and hyperparameter optimization})~\cite{autoweka} es
\[
(j^{*},\lambda^{*})\in\argmin_{j\in\{1,\dots,R\},\ \lambda\in\Lambda^{(j)}}\ \frac1k\sum_{i=1}^{k}\mathcal{L}\bigl(A^{(j)}_{\lambda},D^{(i)}_{\mathrm{tr}},D^{(i)}_{\mathrm{va}}\bigr),
\]
donde $\mathcal{L}$ es la pérdida de validación del modelo entrenado con $D^{(i)}_{\mathrm{tr}}$.
\end{definicion}

El espacio de la \cref{def:cash} es \emph{condicional}: la profundidad de un árbol solo existe si se elige un árbol. Se representa como un único espacio jerárquico cuya raíz es el hiperparámetro «algoritmo», de modo que métodos como la optimización bayesiana con bosques aleatorios (SMAC) o con estimadores de densidad en árbol (TPE) lo tratan sin cambios. Si además se eligen el preprocesado, la selección de variables y la combinación de modelos, la configuración es un programa completo; TPOT, por ejemplo, busca esas canalizaciones con programación genética, como en el \cref{tema:1}~\cite{automl}.

\begin{observacion}[Por qué la búsqueda aleatoria es una línea base seria]\label{obs:aleatoria}
Si la pérdida depende sobre todo de unos pocos hiperparámetros, una rejilla desperdicia evaluaciones repitiendo los mismos valores de los importantes, mientras que $n$ puntos aleatorios prueban $n$ valores distintos de cada uno~\cite{bergstra}. Además, la probabilidad de que al menos uno de $n$ puntos aleatorios caiga en la fracción $q$ mejor del espacio es $1-(1-q)^{n}$, independiente de la dimensión: con $q=0{,}05$ y $n=60$ vale $0{,}954$ (módulo~25, tema~6).
\end{observacion}

\section{Búsqueda de arquitecturas neuronales}

Elsken, Metzen y Hutter organizan la búsqueda de arquitecturas (NAS) en tres dimensiones~\cite{elsken}: el \emph{espacio de búsqueda} (qué arquitecturas pueden representarse), la \emph{estrategia de búsqueda} (cómo se explora) y la \emph{estimación del rendimiento} (cómo se evalúa una arquitectura sin entrenarla del todo).

\begin{definicion}[Espacio de celdas]\label{def:celdas}
Una \emph{celda} es un grafo dirigido acíclico con dos nodos de entrada (las salidas de las dos celdas anteriores) y $B$ nodos intermedios. El nodo intermedio $b$ elige dos entradas entre los $b+1$ nodos anteriores y aplica a cada una una operación de un conjunto $\mathcal{O}$ (convoluciones, agrupamientos, identidad\dots); su salida es la suma de ambas. La red se construye apilando copias de una celda normal y de una celda de reducción.
\end{definicion}

Si las dos elecciones de cada nodo son ordenadas y pueden repetirse, el número de celdas distintas es $\prod_{b=1}^{B}\bigl((b+1)\abs{\mathcal{O}}\bigr)^{2}$: con $B=5$ y $\abs{\mathcal{O}}=8$, unas $5{,}6\cdot10^{14}$. La celda reduce el espacio respecto a una red arbitraria, pero sigue siendo inabarcable por enumeración.

\emph{Estrategias de búsqueda.} Zoph y Le entrenaron con refuerzo una red recurrente que genera descripciones de arquitecturas y obtuvieron un $3{,}65\,\%$ de error en CIFAR-10~\cite{zoph}. Real \emph{et al.} usaron evolución por torneo con \emph{envejecimiento} (se elimina el individuo más antiguo, no el peor), y en una comparación controlada con el mismo hardware la evolución llegó antes a buenos resultados que el refuerzo~\cite{real}. Liu, Simonyan y Yang hicieron el espacio continuo:

\begin{definicion}[Relajación continua de DARTS]\label{def:darts}
En DARTS~\cite{darts}, cada arista $(i,j)$ de la celda aplica una \emph{mezcla} de todas las operaciones,
\[
\bar o^{(i,j)}(x)=\sum_{o\in\mathcal{O}}\frac{\exp\bigl(\alpha^{(i,j)}_{o}\bigr)}{\sum_{o'\in\mathcal{O}}\exp\bigl(\alpha^{(i,j)}_{o'}\bigr)}\,o(x),
\]
y los parámetros de arquitectura $\alpha$ se optimizan por descenso de gradiente en un problema binivel: $\min_{\alpha}\mathcal{L}_{\mathrm{va}}\bigl(w^{*}(\alpha),\alpha\bigr)$ con $w^{*}(\alpha)=\argmin_{w}\mathcal{L}_{\mathrm{tr}}(w,\alpha)$. Al terminar, cada arista conserva la operación de mayor $\alpha$.
\end{definicion}

\emph{Estimación del rendimiento.} Entrenar cada candidata hasta el final es lo que hace prohibitiva la búsqueda. ENAS comparte los pesos entre todas las arquitecturas del espacio (cada candidata es un subgrafo de un modelo grande) y lo abarata en unas mil veces en horas de GPU respecto a la búsqueda original~\cite{enas}. Li y Talwalkar mostraron que la búsqueda aleatoria con parada temprana iguala a ENAS en PTB y CIFAR-10, y con pesos compartidos lo supera, y señalaron problemas de reproducibilidad en la literatura~\cite{lt}. La lección práctica es la del \cref{tema:1}: compárese siempre con una búsqueda aleatoria del mismo presupuesto.

\section{Multifidelidad y coste}

\begin{definicion}[Fidelidad]\label{def:fidelidad}
Una \emph{fidelidad} es un recurso $r\in[r_{\min},r_{\max}]$ que controla el coste y la precisión de la evaluación de una configuración $\lambda$: épocas de entrenamiento, fracción de los datos, resolución de las imágenes o número de árboles. Se escribe $f(\lambda,r)$ para la pérdida de validación con presupuesto $r$; la de interés es $f(\lambda,r_{\max})$.
\end{definicion}

\begin{definicion}[Halving sucesivo]\label{def:sh}
Dados $N$ configuraciones, un factor $\eta\ge2$ y un presupuesto mínimo $r_{\min}$, el \emph{halving sucesivo}~\cite{sh} evalúa todas con $r_0=r_{\min}$, conserva la fracción $1/\eta$ mejor, multiplica el presupuesto por $\eta$ y repite: en el escalón $i$ quedan $n_i=N/\eta^{i}$ configuraciones con presupuesto $r_i=r_{\min}\eta^{i}$, hasta $r_K=r_{\max}$.
\end{definicion}

\begin{proposicion}[Coste del halving sucesivo]\label{prop:coste-sh}
Si $N$ es múltiplo de $\eta^{K}$ y el coste de una evaluación es proporcional a su presupuesto, cada escalón cuesta $N r_{\min}$ y el total es $(K+1)Nr_{\min}$, frente a $Nr_{\max}=N r_{\min}\eta^{K}$ de evaluar todas con el presupuesto completo. El cociente es $(K+1)/\eta^{K}$.
\end{proposicion}

\begin{proposicion}[Cuándo acierta el halving sucesivo]\label{prop:acierto-sh}
Sea $\lambda^{\dagger}$ la mejor configuración a fidelidad completa. Si en cada escalón $i<K$ la ordenación por $f(\cdot,r_i)$ coloca a $\lambda^{\dagger}$ entre las $n_{i+1}$ primeras de las que quedan, el halving sucesivo devuelve $\lambda^{\dagger}$. Si las curvas de aprendizaje se cruzan (una configuración peor al principio es mejor al final), puede descartarla.
\end{proposicion}

La hipótesis de la \cref{prop:acierto-sh} se puede comprobar en un piloto con la correlación de rangos de Spearman entre las pérdidas a baja y a alta fidelidad. Hyperband~\cite{hyperband} protege frente a una mala elección de $r_{\min}$ ejecutando varias rondas de halving sucesivo con distinta agresividad (módulo~25, tema~6). En la práctica, el coste no es exactamente proporcional al presupuesto: cargar los datos, compilar el modelo o arrancar una máquina son costes fijos que reducen el ahorro de la \cref{prop:coste-sh}.

\section{Caso de estudio: AutoML del clasificador de averías}

\addtocounter{definicion}{1}%

\begin{table}[h]
\centering\small
\caption{Tecnologías del tema~3.}\label{tab:tec3}
\begin{tabularx}{\textwidth}{@{}>{\raggedright\arraybackslash}p{2.5cm}>{\raggedright\arraybackslash}X>{\raggedright\arraybackslash}p{2.1cm}>{\raggedright\arraybackslash}X@{}}
\toprule
Herramienta & Para qué sirve & Licencia o tipo & Cuándo elegirla \\
\midrule
Optuna~\cite{optuna} & optimización de hiperparámetros con espacios definidos en el propio código (TPE y otros) y poda con halving sucesivo o Hyperband & abierta (MIT) & ajuste de cualquier modelo en Python con poco código \\
Ray Tune~\cite{ray} & ajuste distribuido en un clúster con planificadores de multifidelidad & abierta (Apache~2.0) & muchas pruebas en paralelo sobre varias máquinas \\
SMAC3 & configuración algorítmica con bosques aleatorios como sustituto & abierta (BSD-3) & espacios condicionales, configurar resolvedores \\
auto-sklearn, AutoGluon & AutoML completo para datos tabulares (selección, ajuste, conjuntos) & abiertas (BSD-3; Apache~2.0) & primera línea base fuerte en datos tabulares \\
NNI & búsqueda de arquitecturas y de hiperparámetros & abierta (MIT) & NAS con estrategias variadas sobre PyTorch \\
\bottomrule
\end{tabularx}
\end{table}

\addtocounter{definicion}{1}%


% =====================================================================
\chapter{Del prototipo a la producción: arquitectura, MLOps y evaluación}
\label{tema:4}
% =====================================================================

Un cuaderno que predice averías y propone un calendario es un prototipo; un sistema que lo hace cada día durante años, con datos que cambian, usuarios que dependen de él y auditores que preguntan por qué decidió algo, es otra cosa. Este tema describe esa otra cosa. La arquitectura y el ciclo de vida siguen el curso CS329S de Stanford~\cite{cs329s} y el libro de Huyen~\cite{huyen}; la deuda técnica, a Sculley \emph{et al.}~\cite{sculley}; la lista de pruebas, a Breck \emph{et al.}~\cite{breck}; la evaluación contrafactual, a Dudík, Langford y Li~\cite{dudik}; la fiabilidad, al capítulo~4 «Optimum Maintenance Policies» de Barlow y Proschan~\cite{bp} y al curso 1.010 del MIT~\cite{mit1010}. Completa la arquitectura del módulo~24 (tema~1, §1.2) y la puesta en planta del módulo~26 (tema~6), que trató el control por capas, los modos de operación, las alarmas y la ciberseguridad.

\section{Arquitectura de un optimizador en producción}

\addtocounter{figure}{1}%

La \cref{fig:produccion} separa responsabilidades que en el prototipo estaban mezcladas.
\begin{enumerate}[label=\arabic*.]
  \item \emph{Datos}: la ingesta valida un \emph{esquema} (tipos, rangos, unidades, frecuencias, ausencias) antes de aceptar nada; los datos y las características se \emph{versionan}, de modo que se pueda reconstruir exactamente con qué datos se entrenó cada modelo.
  \item \emph{Características}: se calculan con el mismo código al entrenar y al servir, y respetando el tiempo: cada fila de entrenamiento solo puede usar información disponible en su instante.
  \item \emph{Modelos}: se entrenan con una canalización reproducible y se guardan en un \emph{registro} con sus métricas, sus datos de origen y su estado (candidato, en sombra, en producción, retirado).
  \item \emph{Servicio}: una API con un contrato explícito (entradas, salidas, unidades, versión del modelo, latencia máxima) desacopla a los consumidores del modelo concreto.
  \item \emph{Optimizador}: combina las predicciones con el modelo de decisión y un resolvedor (módulo~24); sus restricciones duras no dependen del modelo aprendido.
  \item \emph{Interfaz humana}: presenta la propuesta con sus motivos, permite aprobarla, modificarla o anularla, y registra la decisión final; es el lugar donde se ejerce la supervisión que exige la regulación (\cref{tema:5}).
  \item \emph{Monitorización}: vigila los datos, las predicciones, las decisiones, los indicadores de negocio y la salud del servicio, y dispara acciones definidas de antemano.
\end{enumerate}

\begin{definicion}[Sesgo entre entrenamiento y servicio]\label{def:skew}
Hay \emph{sesgo entre entrenamiento y servicio} cuando la distribución de las entradas que el modelo recibe en producción difiere de la de entrenamiento por causas del propio sistema y no del mundo: características calculadas con código distinto, unidades distintas, datos que en entrenamiento estaban disponibles y en servicio llegan tarde, o filas de entrenamiento que contienen información del futuro.
\end{definicion}

\section{Deuda técnica en sistemas con aprendizaje automático}

Sculley \emph{et al.}~\cite{sculley} observan que en un sistema real el código del modelo es una pequeña parte del total y enumeran riesgos propios de estos sistemas, que generan \emph{deuda técnica}: costes futuros que se contraen al tomar atajos hoy.

\begin{definicion}[Factores de riesgo de Sculley \emph{et al.}]\label{def:deuda}
\begin{enumerate}[label=(\alph*)]
  \item \emph{Erosión de fronteras y entrelazamiento}: en un modelo, cambiar cualquier entrada cambia todo lo demás (principio CACE, \emph{changing anything changes everything}); no hay módulos independientes.
  \item \emph{Bucles de realimentación ocultos}: las decisiones del sistema cambian los datos con que se reentrena.
  \item \emph{Consumidores no declarados}: otros sistemas usan las salidas del modelo sin que nadie lo sepa, y se rompen al cambiarlo.
  \item \emph{Dependencias de datos}: entradas inestables (su significado cambia) o infrautilizadas (aportan poco y añaden riesgo).
  \item \emph{Antipatrones de sistema}: código pegamento, junglas de canalizaciones, rutas experimentales muertas.
  \item \emph{Deuda de configuración}: miles de parámetros de configuración sin revisión ni pruebas.
  \item \emph{Cambios en el mundo exterior}: umbrales fijos que dejan de valer, correlaciones que desaparecen.
\end{enumerate}
\end{definicion}

\begin{observacion}[El bucle oculto del mantenimiento predictivo]\label{obs:bucle}
Si el sistema adelanta las intervenciones de las máquinas con riesgo alto, esas máquinas fallarán menos, y los datos futuros mostrarán pocos fallos asociados a vibraciones altas, porque se intervino antes. Un modelo reentrenado ingenuamente con esos datos concluirá que la vibración alta es menos peligrosa y relajará las intervenciones. Es el bucle de realimentación degenerado de la clase~10 de CS329S~\cite{cs329s}. Remedios: registrar cada decisión con su probabilidad (\cref{sec:contrafactual}), tratar las intervenciones como censura en el modelo de supervivencia y mantener una pequeña fracción de decisiones exploratorias cuando sea seguro.
\end{observacion}

Breck \emph{et al.}~\cite{breck} convierten estos riesgos en una lista de 28 pruebas agrupadas en cuatro bloques: datos y características, desarrollo del modelo, infraestructura y monitorización; cada prueba manual puntúa medio punto y cada prueba automatizada, uno. Es una buena lista de comprobación antes de desplegar.

\section{Deriva y reentrenamiento}

\begin{definicion}[Tipos de cambio de distribución]\label{def:deriva}
Sea $X$ la entrada y $Y$ la salida, con $P(X,Y)=P(Y\mid X)P(X)=P(X\mid Y)P(Y)$. Siguiendo la clase~10 de CS329S~\cite{cs329s}, hay \emph{cambio de covariables} si cambia $P(X)$ y no $P(Y\mid X)$; \emph{cambio de etiquetas} si cambia $P(Y)$ y no $P(X\mid Y)$; y \emph{deriva de concepto} si cambia $P(Y\mid X)$. Además pueden cambiar las propias características (unidades, definición) o el esquema de las etiquetas.
\end{definicion}

\begin{definicion}[Índice de estabilidad poblacional]\label{def:psi}
Sea una variable discretizada en $B$ intervalos, con proporciones $p_b$ en la referencia (entrenamiento) y $q_b$ en una ventana reciente, todas positivas. El \emph{índice de estabilidad poblacional} es
\[
\PSI(p,q)=\sum_{b=1}^{B}(q_b-p_b)\ln\frac{q_b}{p_b}.
\]
\end{definicion}

\begin{proposicion}[El PSI es la divergencia simetrizada]\label{prop:psi}
$\PSI(p,q)=\KL(p\,\|\,q)+\KL(q\,\|\,p)\ge 0$, con igualdad si y solo si $p=q$.
\end{proposicion}

\addtocounter{definicion}{1}%

Las métricas de exactitud son la señal más directa de degradación, pero sus etiquetas llegan con retraso (en el caso de estudio, 7 días). Mientras tanto se vigilan las entradas (PSI, contraste de Kolmogorov--Smirnov por variable) y la distribución de las predicciones, más barata por ser de baja dimensión~\cite{cs329s}. Muchas alarmas de deriva resultan ser errores internos: valores ausentes, cambios de unidades, canalizaciones rotas.

\begin{proposicion}[Periodo óptimo de reentrenamiento]\label{prop:reentreno}
Supóngase que, tras reentrenar, la pérdida diaria por degradación del modelo crece linealmente, $at$ (€/día) con $a>0$, y que cada reentrenamiento (cómputo, validación, despliegue) cuesta $c_r$. Con un periodo fijo $T$, el coste medio por día es $c_r/T+aT/2$, que se minimiza en $T^{*}=\sqrt{2c_r/a}$ con coste $\sqrt{2ac_r}$.
\end{proposicion}

Es la fórmula del lote económico aplicada al modelo. En la práctica se combina con disparadores: se reentrena antes si la monitorización detecta una deriva o una caída de las métricas, y nunca se despliega un modelo reentrenado sin pasar las mismas validaciones que el primero.

\section{Evaluar decisiones sin experimentar}\label{sec:contrafactual}

Para saber si una política nueva es mejor que la actual hay tres caminos: un experimento en la planta (\cref{tema:5}), un simulador (el gemelo digital del módulo~26, tema~2) o los registros históricos. Los registros solo muestran el resultado de la decisión que se tomó; el de las demás es \emph{contrafactual}.

\begin{proposicion}[Estimador por ponderación inversa de la propensión]\label{prop:ips}
Sean registros $(x_i,a_i,r_i)$, $i=1,\dots,n$, independientes, donde el contexto $x_i$ sigue una distribución fija, la acción $a_i$ se eligió con una política de registro conocida $\pi_0(a\mid x)$ y $r_i$ es la recompensa observada. Sea $\pi$ una política nueva con $\pi(a\mid x)>0\Rightarrow\pi_0(a\mid x)>0$. Entonces
\[
\hat V_{\mathrm{IPS}}(\pi)=\frac1n\sum_{i=1}^{n}\frac{\pi(a_i\mid x_i)}{\pi_0(a_i\mid x_i)}\,r_i
\]
es un estimador insesgado del valor $V(\pi)=\E_{x}\E_{a\sim\pi(\cdot\mid x)}\bigl[\E[r\mid x,a]\bigr]$.
\end{proposicion}

La varianza crece con los cocientes $\pi/\pi_0$: si la política antigua casi nunca tomaba la acción que la nueva prefiere, la estimación es inútil. El estimador doblemente robusto de Dudík \emph{et al.}~\cite{dudik} combina un modelo de la recompensa con la ponderación y es consistente si cualquiera de los dos es correcto. La consecuencia de diseño es sencilla: \emph{registrar siempre la acción, el contexto y la probabilidad con que se eligió}; un sistema determinista que no explora nunca permite evaluar contrafactuales con sus propios datos.

\section{Fiabilidad y mantenimiento}

\begin{definicion}[Fiabilidad y tasa de fallo]\label{def:fiabilidad}
Sea $T\ge0$ la vida útil de un componente, con función de distribución $F$ y densidad $f$. La \emph{fiabilidad} es $R(t)=1-F(t)$ y la \emph{tasa de fallo} es $\lambda(t)=f(t)/R(t)$: $\lambda(t)\,dt$ es la probabilidad de fallar en $(t,t+dt)$ habiendo sobrevivido hasta $t$. Se cumple $R(t)=\exp\bigl(-\int_0^t\lambda(u)\,du\bigr)$~\cite{mit1010}. La vida es de tasa creciente (IFR) si $\lambda$ es no decreciente y de tasa decreciente (DFR) si es no creciente.
\end{definicion}

\begin{definicion}[Distribución de Weibull]\label{def:weibull}
$T$ sigue una Weibull de forma $\beta>0$ y escala $\theta>0$ si $R(t)=\exp\bigl(-(t/\theta)^{\beta}\bigr)$. Su tasa de fallo es $\lambda(t)=\frac{\beta}{\theta}\bigl(\frac{t}{\theta}\bigr)^{\beta-1}$: decreciente si $\beta<1$ (rodaje), constante si $\beta=1$ (exponencial, sin memoria) y creciente si $\beta>1$ (desgaste). Su media es $\theta\,\Gamma(1+1/\beta)$.
\end{definicion}

\begin{observacion}[Censura]\label{obs:censura}
En los históricos de mantenimiento, muchas vidas no terminan en fallo: la máquina se intervino antes. Esas observaciones están \emph{censuradas por la derecha}: solo se sabe que $T>t_c$. La verosimilitud correcta es $\prod_{\text{fallos}}f(t_i)\prod_{\text{censuras}}R(t_c)$. Descartar las censuras, o tratarlas como fallos, sesga la estimación porque elimina precisamente las vidas largas.
\end{observacion}

\begin{proposicion}[Política de sustitución por edad]\label{prop:edad}
Bajo la política «intervenir preventivamente al alcanzar la edad $T$, o correctivamente al fallar antes», con costes $c_p<c_f$ y renovación completa tras cada intervención, el coste medio a largo plazo por unidad de tiempo es
\[
C(T)=\frac{c_pR(T)+c_fF(T)}{\int_0^{T}R(u)\,du}.
\]
Sea $G(T)=\lambda(T)\int_0^TR(u)\,du-F(T)$. Entonces $C'(T)$ tiene el signo de $G(T)-c_p/(c_f-c_p)$, y $G'(T)=\lambda'(T)\int_0^TR(u)\,du$. En consecuencia:
\begin{enumerate}
  \item si la vida es DFR (en particular, Weibull con $\beta\le1$), $C$ es decreciente y lo óptimo es no intervenir preventivamente ($T^{*}=\infty$);
  \item si $\lambda$ es creciente y $\lambda(T)\to\infty$ (Weibull con $\beta>1$), existe un único $T^{*}$ con $G(T^{*})=c_p/(c_f-c_p)$, y en él $C(T^{*})=(c_f-c_p)\lambda(T^{*})$.
\end{enumerate}
\end{proposicion}

La política por edad usa solo el tiempo. El mantenimiento \emph{predictivo} usa además el estado medido de cada máquina: un clasificador estima la probabilidad de fallo en el horizonte y la decisión se toma con esa probabilidad, como en la decisión con predicciones del módulo~24 (tema~6, §6.1). Si la probabilidad de fallo en los próximos 7 días es $p_7$, intervenir hoy compensa frente a no intervenir en la semana cuando $c_p<c_fp_7$, es decir, cuando $p_7>c_p/c_f$; las restricciones de capacidad y de calendario convierten esa regla en un problema de asignación.

\section{Caso de estudio: mantenimiento predictivo de una línea de envasado}

\addtocounter{definicion}{1}%

\begin{table}[h]
\centering\small
\caption{Costes de cada opción en un lunes con $p_7=(0{,}42;\ 0{,}18;\ 0{,}07;\ 0{,}31;\ 0{,}03;\ 0{,}12)$ y horizonte de una semana (en €). En negrita, la solución óptima con un técnico por ventana.}\label{tab:plan}
\setlength{\tabcolsep}{4pt}\footnotesize
\begin{tabular}{@{}lcccccc@{}}
\toprule
 & llenadora & taponadora & etiquetadora & encajonadora & paletizadora & enfardadora \\
\midrule
$\lambda_i$ (1/día) & 0{,}0778 & 0{,}0284 & 0{,}0104 & 0{,}0530 & 0{,}0044 & 0{,}0183 \\
lunes ($d=0$) & \textbf{1200} & 1200 & 1200 & 1200 & 1200 & 1200 \\
miércoles ($d=2$) & 2497 & 1696 & 1385 & \textbf{2105} & 1278 & 1523 \\
viernes ($d=4$) & 3607 & 2165 & 1566 & 2920 & 1355 & 1834 \\
no intervenir & 3780 & \textbf{1620} & \textbf{630} & 2790 & \textbf{270} & \textbf{1080} \\
\bottomrule
\end{tabular}
\end{table}

Con un horizonte de una semana (\cref{tab:plan}), la llenadora, la taponadora y la encajonadora preferirían el lunes, pero solo hay un técnico. El ahorro respecto a no intervenir es $3780-1200=2580$ para la llenadora el lunes y $2790-2105=685$ para la encajonadora el miércoles; la alternativa (encajonadora el lunes y llenadora el miércoles) ahorra $1590+1283=2873<3265$. El coste óptimo es $10\,170-3265=6905$~€. Con el horizonte de dos semanas, el modelo añade además la taponadora el viernes, porque no intervenirla en catorce días es más arriesgado que en siete.

\begin{table}[h]
\centering\small
\caption{Un año de operación de la línea en el gemelo, con los mismos números aleatorios para todas las políticas (módulo~26, tema~5, §5.1).}\label{tab:politicas}
\setlength{\tabcolsep}{5pt}
\begin{tabular}{@{}lrrrr@{}}
\toprule
Política & Coste (€/año) & Averías & Preventivos & Fuera de ventana \\
\midrule
A. Solo correctivo & 360\,000 & 40 & 0 & --- \\
H. Actual: cada 70 días & 327\,400 & 35 & 7 & 5 \\
B. Por edad, $T^{*}=30$ días & 195\,200 & 10 & 65 & 34 \\
C. Umbral $p_7\ge0{,}10$, sin calendario & 173\,200 & 6 & 70 & 44 \\
D. Calendario con riesgo (PLEM) & 141\,000 & 7 & 65 & 0 \\
\bottomrule
\end{tabular}
\end{table}

El \cref{tab:politicas} resume la simulación del cuaderno. El umbral de la política C ($0{,}10$, cercano a $c_p/c_f$) se eligió en un año de validación distinto. La política C, además de intervenir fuera de ventana, viola la capacidad del técnico cinco días. La política D respeta todas las restricciones y es la más barata. Que la mejora frente a B sea real, y no suerte de este año, es la pregunta del \cref{tema:5}.

\begin{table}[h]
\centering\small
\caption{Tecnologías del tema~4.}\label{tab:tec4}
\begin{tabularx}{\textwidth}{@{}>{\raggedright\arraybackslash}p{2.5cm}>{\raggedright\arraybackslash}X>{\raggedright\arraybackslash}p{2.1cm}>{\raggedright\arraybackslash}X@{}}
\toprule
Herramienta & Para qué sirve & Licencia o tipo & Cuándo elegirla \\
\midrule
DVC~\cite{dvc} & versionar datos y modelos junto al código, con canalizaciones reproducibles & abierta (Apache~2.0) & equipos pequeños que ya usan git \\
MLflow~\cite{mlflow} & seguimiento de experimentos, registro de modelos y empaquetado & abierta (Apache~2.0) & trazabilidad de qué modelo está en producción \\
Great Expectations, Feast & validación de esquemas y almacén de características & abiertas (Apache~2.0) & evitar el sesgo de la \cref{def:skew} \\
Evidently~\cite{evidently} & informes y monitorización de deriva de datos y predicciones & abierta (Apache~2.0) & vigilancia periódica con pocos recursos \\
FastAPI~\cite{fastapi} & servicio de predicción y optimización como API HTTP con esquema & abierta (MIT) & exponer el modelo y el resolvedor a otros sistemas \\
Kubernetes~\cite{k8s}, Ray Serve~\cite{ray} & despliegue en contenedores, escalado, versiones en paralelo & abiertas (Apache~2.0) & varios servicios con requisitos de disponibilidad \\
lifelines, scikit-survival & modelos de supervivencia con censura & abiertas (MIT; GPL-3.0) & estimar fiabilidad y riesgo en función del tiempo \\
\bottomrule
\end{tabularx}
\end{table}

\addtocounter{definicion}{1}%


% =====================================================================
\chapter{Cuantificar el valor y el riesgo}
\label{tema:5}
% =====================================================================

Un optimizador se justifica si su valor esperado supera a su coste y si sus riesgos son aceptables y están controlados. Este tema da las herramientas para afirmarlo con números: intervalos de confianza para la mejora con diseños pareados (capítulo~2 de Montgomery~\cite{montgomery}) y con \emph{bootstrap} (lectura~24 del curso 18.05 del MIT~\cite{mit1805}), tamaño de un piloto, valor actual neto (sesiones «Present Value Relations» y «Capital Budgeting» del curso 15.401 del MIT~\cite{mit15401}), medidas de riesgo de cola (Rockafellar y Uryasev~\cite{ru}), análisis de modos de fallo (IEC~60812~\cite{iec60812}) y el marco del Reglamento europeo de inteligencia artificial~\cite{aiact,omnibus}. Se apoya en el módulo~0.9 (temas~4 y~5), en el contraste del ahorro del módulo~26 (tema~6, ejemplo «Contraste del ahorro») y en los experimentos en línea de Kohavi, Tang y Xu~\cite{kohavi}.

\section{Coste y beneficio}

\begin{definicion}[Valor actual neto]\label{def:van}
Un proyecto con inversión inicial $I$ y flujos netos $B_t-O_t$ (beneficios menos costes de operación) al final de cada año $t=1,\dots,N$, con tipo de descuento anual $\rho>0$, tiene \emph{valor actual neto}
\[
\VAN=-I+\sum_{t=1}^{N}\frac{B_t-O_t}{(1+\rho)^{t}}.
\]
La \emph{tasa interna de rentabilidad} (TIR) es el $\rho$ que anula el VAN, y el \emph{plazo de recuperación}, el primer año en que los flujos acumulados sin descontar cubren $I$.
\end{definicion}

\begin{proposicion}[Anualidad]\label{prop:anualidad}
Si $B_t-O_t=b$ es constante, $\VAN=-I+b\,a_{N,\rho}$ con $a_{N,\rho}=\sum_{t=1}^{N}(1+\rho)^{-t}=\dfrac{1-(1+\rho)^{-N}}{\rho}$.
\end{proposicion}

Los beneficios de un optimizador se miden \emph{respecto a una alternativa}, y la elección de la alternativa decide el resultado. La alternativa honrada no es la práctica actual, sino la mejor solución sencilla que se adoptaría sin IA. En el mantenimiento del \cref{tema:4}, pasar de intervenir cada 70 días a hacerlo cada 30 (política~B) no requiere sensores ni modelos de clasificación: basta un análisis de fiabilidad. El valor atribuible al optimizador con IA es la diferencia entre la política~D y la~B.

\section{Intervalos de confianza para la mejora}

\begin{proposicion}[Ventaja del diseño pareado]\label{prop:pareado}
Sean $Y_A$ e $Y_B$ los costes de dos políticas en el mismo periodo, con varianzas $\sigma_A^{2}$, $\sigma_B^{2}$ y correlación $r$. La diferencia pareada $D=Y_A-Y_B$ tiene varianza $\sigma_A^{2}+\sigma_B^{2}-2r\sigma_A\sigma_B$, menor que la de la diferencia de dos medidas independientes, $\sigma_A^{2}+\sigma_B^{2}$, si y solo si $r>0$.
\end{proposicion}

Comparar dos políticas en los mismos días, o en el gemelo con los mismos números aleatorios (módulo~26, tema~5, §5.1), crea esa correlación positiva: la carga, el desgaste y el azar comunes se cancelan. Con $n$ diferencias pareadas de media $\bar d$ y desviación típica $s_d$, el intervalo de confianza de nivel $1-\alpha$ para la mejora media es $\bar d\pm t_{n-1,1-\alpha/2}\,s_d/\sqrt n$ (módulo~0.9, tema~4, §4.4).

\begin{definicion}[Intervalos \emph{bootstrap}]\label{def:bootstrap}
Sean $d_1,\dots,d_n$ las diferencias observadas y $\bar d$ su media. Se generan $B$ remuestras de tamaño $n$ con reemplazamiento, con medias $\bar d^{*}_1,\dots,\bar d^{*}_B$. El intervalo \emph{percentil} de nivel $1-\alpha$ es $[\bar d^{*}_{(\alpha/2)},\bar d^{*}_{(1-\alpha/2)}]$, con los cuantiles empíricos de las medias remuestreadas. El intervalo \emph{básico} usa $\delta^{*}=\bar d^{*}-\bar d$: $[\bar d-\delta^{*}_{(1-\alpha/2)},\ \bar d-\delta^{*}_{(\alpha/2)}]$~\cite{mit1805}.
\end{definicion}

El \emph{bootstrap} no exige normalidad, lo que importa aquí: el coste semanal de una línea es casi siempre pequeño y, de vez en cuando, salta en 9000~€ por una avería. Ambos métodos suponen observaciones independientes; si las semanas están autocorrelacionadas hay que remuestrear por bloques.

\begin{proposicion}[Tamaño de un piloto]\label{prop:tamano}
Para detectar una mejora media $\delta>0$ con un contraste bilateral de nivel $\alpha$ y potencia $1-\beta$ con diferencias pareadas de desviación típica $\sigma_D$, hacen falta aproximadamente
\[
n\approx\Bigl(\frac{(z_{1-\alpha/2}+z_{1-\beta})\,\sigma_D}{\delta}\Bigr)^{2}
\]
pares (módulo~0.9, tema~5, §5.4).
\end{proposicion}

\section{Caso de estudio: ¿cuánto vale el optimizador de mantenimiento?}

\addtocounter{definicion}{1}%

\section{Robustez y riesgo de cola}

La media no lo es todo: a una planta le preocupa el año malo. Para una pérdida $L$ con distribución continua, el \emph{valor en riesgo} al nivel $\alpha$ es el cuantil $\mathrm{VaR}_\alpha(L)=F_L^{-1}(\alpha)$, y el \emph{valor en riesgo condicional} es la media de la cola, $\mathrm{CVaR}_\alpha(L)=\E[L\mid L\ge\mathrm{VaR}_\alpha(L)]$.

\begin{proposicion}[Fórmula de Rockafellar y Uryasev]\label{prop:cvar}
Si $L$ tiene densidad y esperanza finita,
\[
\mathrm{CVaR}_\alpha(L)=\min_{\tau\in\R}\Bigl\{\tau+\frac{1}{1-\alpha}\E\bigl[(L-\tau)^{+}\bigr]\Bigr\},
\]
y el mínimo se alcanza en $\tau=\mathrm{VaR}_\alpha(L)$.
\end{proposicion}

La fórmula convierte la minimización del CVaR de una decisión en un programa lineal cuando $L$ se representa con escenarios: basta añadir $\tau$ y una variable $u_s\ge L_s-\tau$, $u_s\ge0$ por escenario~\cite{ru}. En los 20 años simulados del caso, el coste anual de la política~D tiene media $159\,300$~€ y CVaR al 90\,\% (media de los dos peores años) de $186\,300$~€; la política~B, $193\,670$ y $216\,200$~€. D es mejor también en la cola. Las \emph{pruebas de estrés} completan el análisis: repetir la evaluación con un coste de avería un 50\,\% mayor, con un técnico de baja, con un sensor que se congela en su último valor o con una semana de datos perdidos.

\section{Seguridad, supervisión y modos de fallo}

Los principios del módulo~26 (tema~6) se aplican a cualquier optimizador. (i) Las \emph{restricciones duras} (capacidad de los técnicos, inspecciones obligatorias por ley, bloqueos de seguridad de las máquinas) están en el modelo de decisión y se comprueban otra vez, de forma independiente, antes de emitir una orden. (ii) Hay una \emph{política de respaldo} sencilla (aquí, la política por edad) a la que se vuelve si el modelo o los datos fallan. (iii) El despliegue recorre los modos asesor, supervisado y autónomo con supervisión (módulo~26, tema~6, §6.2), y en cada uno se registra quién decidió qué.

\begin{definicion}[Análisis modal de fallos y efectos]\label{def:amfe}
El AMFE (FMEA en inglés; IEC~60812~\cite{iec60812}) enumera los modos de fallo de cada componente, sus efectos y sus causas, y los prioriza. En la variante clásica se puntúan de 1 a 10 la \emph{gravedad} $S$, la \emph{ocurrencia} $O$ y la \emph{detectabilidad} $Dt$ (10 = indetectable), y se ordena por el \emph{número de prioridad de riesgo} $\mathrm{NPR}=S\cdot O\cdot Dt$; la norma admite otras formas de calcularlo y, en todo caso, los modos de gravedad máxima se tratan aunque su NPR sea bajo.
\end{definicion}

\begin{table}[h]
\centering\small
\caption{AMFE resumido del optimizador de mantenimiento.}\label{tab:amfe}
\setlength{\tabcolsep}{4pt}\footnotesize
\begin{tabularx}{\textwidth}{@{}>{\raggedright\arraybackslash}p{2.8cm}>{\raggedright\arraybackslash}X cccc >{\raggedright\arraybackslash}p{3.0cm}@{}}
\toprule
Modo de fallo & Efecto & $S$ & $O$ & $Dt$ & NPR & Medida \\
\midrule
sensor congelado en su último valor & riesgo infravalorado; avería no prevista & 8 & 3 & 7 & 168 & regla de varianza nula; respaldo por edad \\
deriva tras cambio de producto & probabilidades mal calibradas & 6 & 5 & 4 & 120 & PSI semanal; recalibración \\
resolvedor sin solución a tiempo & sin plan el lunes & 5 & 2 & 2 & 20 & plan de la semana anterior; límite de tiempo \\
confianza excesiva del técnico & se ejecutan planes erróneos & 7 & 4 & 6 & 168 & mostrar motivos; auditoría de anulaciones \\
etiquetas mal registradas en el GMAO & modelo entrenado con ruido & 6 & 4 & 6 & 144 & validación de registros; revisión mensual \\
\bottomrule
\end{tabularx}
\end{table}

\section{El Reglamento europeo de inteligencia artificial}

El Reglamento (UE) 2024/1689~\cite{aiact} entró en vigor el 1 de agosto de 2024 y sigue un enfoque basado en el riesgo: prohíbe algunas prácticas, impone requisitos a los sistemas de \emph{alto riesgo}, establece obligaciones de transparencia para otros y deja el resto sin obligaciones específicas. Lo que sigue es una descripción general de su texto público; la clasificación de un sistema concreto es una cuestión jurídica.

\begin{enumerate}[label=\arabic*.]
  \item \emph{Alto riesgo} (artículo~6). Lo es un sistema que sea componente de seguridad de un producto, o el propio producto, regulado por la legislación de armonización del anexo~I y sujeto a evaluación de la conformidad por terceros; y lo son los sistemas del anexo~III, entre ellos los componentes de seguridad en la gestión y operación de infraestructuras críticas, como el suministro de agua, gas, calefacción o electricidad (anexo~III, punto~2). Un sistema del anexo~III no es de alto riesgo si no influye de forma sustancial en la decisión, por ejemplo si realiza una tarea procedimental limitada o preparatoria (artículo~6, apartado~3), salvo que elabore perfiles de personas.
  \item \emph{Requisitos} (artículos~8 a~15): sistema de gestión de riesgos, gobernanza de los datos, documentación técnica, registro automático de eventos, transparencia hacia quien lo despliega, supervisión humana y precisión, robustez y ciberseguridad. El artículo~14 pide que las personas que supervisan puedan entender las capacidades y los límites del sistema y detectar anomalías, sean conscientes del sesgo de automatización, interpreten correctamente sus salidas, puedan decidir no usarlo o anular sus resultados, e intervenir o detenerlo de forma segura.
  \item \emph{Calendario.} El reglamento se aplica con carácter general desde el 2 de agosto de 2026. La modificación conocida como «ómnibus digital», publicada en el Diario Oficial el 24 de julio de 2026 y en vigor desde el 27 de julio, aplaza las obligaciones de alto riesgo: al 2 de diciembre de 2027 para los sistemas del anexo~III y al 2 de agosto de 2028 para los del anexo~I; según la información publicada, la maquinaria pasa a un régimen sectorial propio~\cite{omnibus,aiact}.
\end{enumerate}

Aplicado al caso: un optimizador del calendario de mantenimiento de una línea de envasado, cuyas órdenes revisa un técnico y que no actúa sobre las funciones de seguridad de las máquinas, no parece encajar en ninguna categoría de alto riesgo. El mismo optimizador aplicado a las bombas de una red de abastecimiento de agua, si actuara como componente de seguridad de su operación, sí podría serlo. En ambos casos, la arquitectura de los temas~4 y~5 (registro de decisiones, supervisión humana con capacidad de anular, política de respaldo, monitorización, documentación) cumple de forma natural buena parte de los requisitos, y es buena práctica adoptarla aunque no sea obligatoria. Si se usan datos de personas (por ejemplo, qué técnico anuló qué orden), se aplica además la normativa de protección de datos.

\begin{table}[h]
\centering\small
\caption{Tecnologías del tema~5.}\label{tab:tec5}
\begin{tabularx}{\textwidth}{@{}>{\raggedright\arraybackslash}p{2.5cm}>{\raggedright\arraybackslash}X>{\raggedright\arraybackslash}p{2.1cm}>{\raggedright\arraybackslash}X@{}}
\toprule
Herramienta & Para qué sirve & Licencia o tipo & Cuándo elegirla \\
\midrule
SciPy (\texttt{stats}) & contrastes $t$, intervalos y \emph{bootstrap} (función \texttt{bootstrap}) & abierta (BSD-3) & análisis de pilotos y de simulaciones \\
statsmodels & modelos con errores autocorrelacionados, potencia y tamaño muestral & abierta (BSD-3) & series semanales, diseños con covariables \\
numpy-financial & VAN, TIR y anualidades & abierta (BSD-3) & valoración económica reproducible \\
MLflow~\cite{mlflow} & registrar cada evaluación con su modelo, datos y resultado & abierta (Apache~2.0) & trazabilidad exigida por la documentación técnica \\
Plantillas de AMFE (hoja de cálculo) & análisis de modos de fallo & --- & revisión de riesgos con operación, mantenimiento y seguridad \\
\bottomrule
\end{tabularx}
\end{table}

\addtocounter{definicion}{1}%


% =====================================================================
\chapter{Guía para construir un optimizador con IA en cualquier ámbito}
\label{tema:6}
% =====================================================================

Este tema reúne la fase~VI en un procedimiento. No es un algoritmo, sino una disciplina de ingeniería: el orden en que conviene tomar las decisiones y los documentos que deben quedar escritos. Sigue la estructura del curso 16.842 «Fundamentals of Systems Engineering» del MIT~\cite{mit16842}, que recorre el modelo en V (partes interesadas y concepto de operación, requisitos con márgenes, arquitectura, exploración del espacio de alternativas, revisión preliminar del diseño), las clases~1 y~2 del curso CS329S de Stanford~\cite{cs329s} y el capítulo~2 de Huyen~\cite{huyen} para el planteamiento de sistemas con aprendizaje automático, y la lista de pruebas de Breck \emph{et al.}~\cite{breck}. Cada paso remite al módulo donde se estudió su contenido.

\section{Los doce pasos}

\begin{enumerate}[label=\textbf{\arabic*.},leftmargin=1.6em]
  \item \textbf{La decisión y sus responsables.} ¿Qué se decide, quién lo decide hoy, con qué frecuencia y con qué horizonte? ¿Quién usará el sistema, quién responde de sus errores y quién puede pararlo? Entregable: concepto de operación de una página.
  \item \textbf{Objetivo, indicadores y restricciones.} Un objetivo escalar, los indicadores que se vigilan sin optimizarlos, las restricciones duras (seguridad, ley, física) y las blandas, cada una con su unidad y su fuente (módulo~24, tema~1). Entregable: requisitos numéricos con márgenes.
  \item \textbf{Línea base y techo de valor.} Coste de la práctica actual, de la mejor alternativa sencilla sin IA y de una política con información perfecta (\cref{prop:techo}). Si el techo no compensa el coste del proyecto, se para aquí.
  \item \textbf{Datos.} Inventario de fuentes, frecuencias, calidad, retrasos, disponibilidad en el instante de decidir, propiedad y protección; versionado desde el primer día (\cref{tema:4}).
  \item \textbf{Formulación y familia de métodos.} El mapa del módulo~24 (tema~1, §1.3): programación matemática, optimización bayesiana para experimentos caros (módulo~25), control y RL para decisiones secuenciales (módulo~26), optimización geométrica (módulo~27), búsqueda de programas cuando lo que se busca es una regla (\cref{tema:1}).
  \item \textbf{Modelos predictivos o sustitutos.} Qué se predice, con qué horizonte y con qué métrica de decisión; validación temporal y selección con presupuesto (\cref{tema:3}).
  \item \textbf{Resolvedor, salvaguardas y respaldo.} Tiempo máximo de cálculo, comprobación independiente de las restricciones duras, política de respaldo.
  \item \textbf{Evaluación fuera de línea.} Sobre históricos (sin fugas), en un gemelo con números aleatorios comunes y, si hay registros con propensiones, con estimadores contrafactuales (\cref{tema:4}).
  \item \textbf{Piloto.} Modo sombra y después experimento con tamaño calculado (\cref{prop:tamano}) y criterio de aceptación escrito antes de empezar.
  \item \textbf{Despliegue.} API con contrato, interfaz que explica y permite anular, registro de decisiones, versiones y vuelta atrás.
  \item \textbf{Monitorización y reentrenamiento.} Entradas, predicciones, decisiones, anulaciones e indicadores; disparadores y periodo de reentrenamiento (\cref{prop:reentreno}).
  \item \textbf{Valor, riesgo, normativa y ciclo de vida.} Intervalos para la mejora, VAN, AMFE, encaje regulatorio, documentación y criterios de retirada (\cref{tema:5}).
\end{enumerate}

\begin{proposicion}[Techo de valor]\label{prop:techo}
Sea $f(x,\theta)$ el coste de la decisión $x\in\mathcal{X}$ cuando el estado incierto es $\theta$, y sea $\hat x$ cualquier regla de decisión que solo use información disponible antes de decidir. Entonces
\[
\E\bigl[f(\hat x,\theta)\bigr]\ \ge\ \E\Bigl[\min_{x\in\mathcal{X}}f(x,\theta)\Bigr].
\]
En consecuencia, ninguna mejora de la predicción puede reducir el coste esperado de la mejor regla sin información por debajo del coste con información perfecta, y la diferencia entre ambos (el valor de la información perfecta del módulo~24, tema~6, definición~6.6) acota lo que puede ganarse prediciendo mejor.
\end{proposicion}

En el caso del mantenimiento, una política que conoce el futuro e interviene cada máquina en la última ventana antes de que fuera a fallar cuesta de media $76\,170$~€/año en el gemelo (sin límite de técnicos), frente a $193\,670$ de la política por edad y $159\,300$ del optimizador. Este solo captura un 29\,\% del margen que deja la alternativa sencilla frente a la información perfecta: hay espacio para mejores sensores o mejores modelos, y el paso~3 permite estimar cuánto valdrían antes de pagarlos.

\section{Plantilla de documento de diseño}

El \cref{tab:plantilla} es la plantilla que se rellena antes de escribir código y se actualiza en cada revisión. Su objetivo no es la burocracia, sino forzar que las decisiones difíciles (qué es éxito, qué no se puede violar, contra qué se compara) se tomen pronto y por escrito, como en la revisión preliminar de diseño de 16.842~\cite{mit16842}.

\begin{table}[h]
\centering\small
\caption{Plantilla del documento de diseño de un optimizador con IA.}\label{tab:plantilla}
\begin{tabularx}{\textwidth}{@{}>{\raggedright\arraybackslash}p{3.6cm}>{\raggedright\arraybackslash}X@{}}
\toprule
Sección & Preguntas que debe responder \\
\midrule
1. Contexto y decisión & ¿Qué se decide, cada cuánto, con qué horizonte? ¿Qué pasa hoy? \\
2. Partes interesadas & ¿Quién decide, quién ejecuta, quién responde, quién puede parar el sistema? \\
3. Objetivo y requisitos & Objetivo, indicadores, restricciones duras y blandas con unidades; requisitos de latencia, disponibilidad y precisión con márgenes. \\
4. Línea base y techo & Coste actual, mejor alternativa sencilla, coste con información perfecta; valor máximo alcanzable. \\
5. Datos & Fuentes, calidad, retrasos, disponibilidad en el instante de decidir, versionado, protección. \\
6. Formulación y método & Modelo matemático; familia elegida y alternativas descartadas con su motivo. \\
7. Modelos y validación & Qué se predice; división temporal; métricas de predicción y de decisión; selección. \\
8. Salvaguardas y respaldo & Comprobaciones independientes; política de respaldo; tiempo máximo de cálculo. \\
9. Evaluación fuera de línea & Históricos, gemelo, contrafactual; resultados con intervalos. \\
10. Piloto & Diseño, tamaño, duración, criterio de aceptación fijado de antemano. \\
11. Despliegue y supervisión & API, interfaz, modos de operación, registro, vuelta atrás. \\
12. Monitorización & Señales, umbrales, acciones, reentrenamiento. \\
13. Valor y riesgo & VAN con incertidumbre, AMFE, pruebas de estrés, encaje normativo. \\
14. Ciclo de vida & Responsable del mantenimiento, revisiones periódicas, criterios de retirada. \\
\bottomrule
\end{tabularx}
\end{table}

\section{Mini-caso A: secuenciación de formatos en una fábrica de bebidas}

\addtocounter{definicion}{1}%

\section{Mini-caso B: presupuesto de geometría en un videojuego}

\addtocounter{definicion}{1}%

\section{Mini-caso C: asignación de trabajos en un centro de datos}

\addtocounter{definicion}{1}%

\begin{table}[h]
\centering\small
\caption{Tecnologías del tema~6: pila de referencia de un optimizador con IA.}\label{tab:tec6}
\begin{tabularx}{\textwidth}{@{}>{\raggedright\arraybackslash}p{2.5cm}>{\raggedright\arraybackslash}X>{\raggedright\arraybackslash}p{2.1cm}>{\raggedright\arraybackslash}X@{}}
\toprule
Capa & Herramientas habituales & Licencia o tipo & Cuándo elegirla \\
\midrule
Datos y versionado & DVC~\cite{dvc}, Great Expectations & abiertas (Apache~2.0) & desde el primer prototipo \\
Modelos & scikit-learn, PyTorch; Optuna~\cite{optuna} para el ajuste & abiertas (BSD-3; BSD-3; MIT) & según el tipo de datos (\cref{tema:3}) \\
Decisión & SciPy/HiGHS, OR-Tools, CVXPY (módulo~24); BoTorch (módulo~25) & abiertas & según la estructura del problema \\
Búsqueda con LLM & OpenEvolve~\cite{openevolve}, DSPy~\cite{dspy} & abiertas (Apache~2.0; MIT) & reglas o cadenas de llamadas con evaluador fiable \\
Servicio & FastAPI~\cite{fastapi}, contenedores, Kubernetes~\cite{k8s} & abiertas (MIT; Apache~2.0) & integración con otros sistemas \\
Seguimiento & MLflow~\cite{mlflow}, Evidently~\cite{evidently} & abiertas (Apache~2.0) & trazabilidad y deriva \\
Escala & Ray~\cite{ray} & abierta (Apache~2.0) & simulaciones y ajustes en paralelo \\
\bottomrule
\end{tabularx}
\end{table}

\addtocounter{definicion}{1}%



\backmatter
% =====================================================================
\begin{thebibliography}{99}
\addcontentsline{toc}{chapter}{Bibliografía}
\markboth{Bibliografía}{}
\small

\bibitem{optuna}
T.~Akiba, S.~Sano, T.~Yanase, T.~Ohta y M.~Koyama.
\newblock Optuna: A next-generation hyperparameter optimization framework.
\newblock En \emph{Proceedings of the 25th ACM SIGKDD Conference on Knowledge Discovery and Data Mining}, 2019. Documentación oficial; licencia MIT.
\newblock \url{https://optuna.org/}

\bibitem{bp}
R.~E. Barlow y F.~Proschan.
\newblock \emph{Mathematical Theory of Reliability}.
\newblock SIAM, colección Classics in Applied Mathematics, 1996 (original: Wiley, 1965). Cap.~4, «Optimum Maintenance Policies».

\bibitem{bergstra}
J.~Bergstra e Y.~Bengio.
\newblock Random search for hyper-parameter optimization.
\newblock \emph{Journal of Machine Learning Research}, 13:281--305, 2012.

\bibitem{bt}
D.~Bertsimas y J.~N. Tsitsiklis.
\newblock \emph{Introduction to Linear Optimization}.
\newblock Athena Scientific, 1997. Cap.~7, «Network flow problems».

\bibitem{breck}
E.~Breck, S.~Cai, E.~Nielsen, M.~Salib y D.~Sculley.
\newblock The ML test score: A rubric for ML production readiness and technical debt reduction.
\newblock En \emph{IEEE International Conference on Big Data}, 2017.

\bibitem{omnibus}
CMS.
\newblock The EU's Digital Omnibus: streamlined AI Act compliance rules are now in force.
\newblock Nota informativa, 2026 (publicación en el Diario Oficial el 24 de julio de 2026; entrada en vigor el 27 de julio de 2026).
\newblock \url{https://cms.law/en/gbr/publication/the-eu-s-digital-omnibus-streamlined-ai-act-compliance-rules-are-now-in-force}

\bibitem{dvc}
Desarrolladores de DVC.
\newblock DVC: Data Version Control. Documentación oficial; licencia Apache~2.0.
\newblock \url{https://dvc.org/doc}

\bibitem{evidently}
Desarrolladores de Evidently.
\newblock Evidently: evaluación y monitorización de modelos de aprendizaje automático. Documentación oficial; licencia Apache~2.0.
\newblock \url{https://docs.evidentlyai.com/}

\bibitem{fastapi}
Desarrolladores de FastAPI.
\newblock FastAPI. Documentación oficial; licencia MIT.
\newblock \url{https://fastapi.tiangolo.com/}

\bibitem{mlflow}
Desarrolladores de MLflow.
\newblock MLflow: plataforma para el ciclo de vida del aprendizaje automático. Documentación oficial; licencia Apache~2.0.
\newblock \url{https://mlflow.org/docs/latest/}

\bibitem{dudik}
M.~Dudík, J.~Langford y L.~Li.
\newblock Doubly robust policy evaluation and learning.
\newblock En \emph{Proceedings of the 28th International Conference on Machine Learning (ICML)}, 2011.
\newblock \url{https://arxiv.org/abs/1103.4601}

\bibitem{elsken}
T.~Elsken, J.~H. Metzen y F.~Hutter.
\newblock Neural architecture search.
\newblock Cap.~3 de~\cite{automl}, 2019.

\bibitem{deepmind}
R.~Evans y J.~Gao.
\newblock DeepMind AI reduces Google data centre cooling bill by 40\,\%.
\newblock Blog de Google DeepMind, 20 de julio de 2016.
\newblock \url{https://deepmind.google/discover/blog/deepmind-ai-reduces-google-data-centre-cooling-bill-by-40/}

\bibitem{ae-blog}
Google DeepMind (equipo de AlphaEvolve).
\newblock AlphaEvolve: A Gemini-powered coding agent for designing advanced algorithms.
\newblock Blog de Google DeepMind, 14 de mayo de 2025.
\newblock \url{https://deepmind.google/discover/blog/alphaevolve-a-gemini-powered-coding-agent-for-designing-advanced-algorithms/}

\bibitem{fs-blog}
A.~Fawzi y B.~Romera-Paredes.
\newblock FunSearch: Making new discoveries in mathematical sciences using large language models.
\newblock Blog de Google DeepMind, 14 de diciembre de 2023 (con una actualización de 2024).
\newblock \url{https://deepmind.google/discover/blog/funsearch-making-new-discoveries-in-mathematical-sciences-using-large-language-models/}

\bibitem{funsearch-code}
Google DeepMind.
\newblock FunSearch: repositorio oficial (implementación del algoritmo evolutivo, cuaderno \texttt{bin\_packing} con los conjuntos OR3 y Weibull 5k y las heurísticas descubiertas). Licencia Apache~2.0.
\newblock \url{https://github.com/google-deepmind/funsearch}

\bibitem{hp}
J.~Herrmann y G.~Pallez.
\newblock An in-depth study of LLM contributions to the bin packing problem.
\newblock Prepublicación, 2025.
\newblock \url{https://arxiv.org/abs/2510.27353}

\bibitem{huyen}
C.~Huyen.
\newblock \emph{Designing Machine Learning Systems}.
\newblock O'Reilly, 2022. Texto asociado al curso CS329S de Stanford.

\bibitem{automl}
F.~Hutter, L.~Kotthoff y J.~Vanschoren (eds.).
\newblock \emph{Automated Machine Learning: Methods, Systems, Challenges}.
\newblock Springer, 2019. Edición de acceso abierto. Cap.~1 «Hyperparameter Optimization» (M.~Feurer y F.~Hutter) y cap.~3 «Neural Architecture Search».
\newblock \url{https://www.automl.org/book/}

\bibitem{iec60812}
IEC.
\newblock IEC 60812:2018. \emph{Failure modes and effects analysis (FMEA and FMECA)}, 3.ª ed.
\newblock International Electrotechnical Commission, 2018. Se usa solo el resumen público.
\newblock \url{https://webstore.iec.ch/en/publication/26359}

\bibitem{sh}
K.~Jamieson y A.~Talwalkar.
\newblock Non-stochastic best arm identification and hyperparameter optimization.
\newblock En \emph{Proceedings of the 19th International Conference on Artificial Intelligence and Statistics (AISTATS)}, 2016.
\newblock \url{https://arxiv.org/abs/1502.07943}

\bibitem{dspy}
O.~Khattab, A.~Singhvi, P.~Maheshwari, Z.~Zhang, K.~Santhanam, S.~Vardhamanan, S.~Haq, A.~Sharma, T.~T. Joshi, H.~Moazam, H.~Miller, M.~Zaharia y C.~Potts.
\newblock DSPy: Compiling declarative language model calls into self-improving pipelines.
\newblock Prepublicación, 2023. Biblioteca con licencia MIT.
\newblock \url{https://arxiv.org/abs/2310.03714}

\bibitem{kw}
M.~J. Kochenderfer y T.~A. Wheeler.
\newblock \emph{Algorithms for Optimization}.
\newblock MIT Press, 2019. Cap.~9 «Population Methods» y cap.~20 «Expression Optimization».
\newblock \url{https://algorithmsbook.com/optimization/}

\bibitem{kohavi}
R.~Kohavi, D.~Tang y Y.~Xu.
\newblock \emph{Trustworthy Online Controlled Experiments: A Practical Guide to A/B Testing}.
\newblock Cambridge University Press, 2020.

\bibitem{k8s}
The Kubernetes Authors.
\newblock Kubernetes. Documentación oficial; licencia Apache~2.0.
\newblock \url{https://kubernetes.io/docs/}

\bibitem{lt}
L.~Li y A.~Talwalkar.
\newblock Random search and reproducibility for neural architecture search.
\newblock En \emph{Proceedings of the 35th Conference on Uncertainty in Artificial Intelligence (UAI)}, 2019.
\newblock \url{https://arxiv.org/abs/1902.07638}

\bibitem{hyperband}
L.~Li, K.~Jamieson, G.~DeSalvo, A.~Rostamizadeh y A.~Talwalkar.
\newblock Hyperband: A novel bandit-based approach to hyperparameter optimization.
\newblock \emph{Journal of Machine Learning Research}, 18(185):1--52, 2018.

\bibitem{darts}
H.~Liu, K.~Simonyan e Y.~Yang.
\newblock DARTS: Differentiable architecture search.
\newblock En \emph{International Conference on Learning Representations (ICLR)}, 2019.
\newblock \url{https://arxiv.org/abs/1806.09055}

\bibitem{mit1010}
MIT OpenCourseWare.
\newblock 1.010: Uncertainty in Engineering.
\newblock Massachusetts Institute of Technology, otoño de 2008. \emph{Example Application}~8, «Old better than new, new better than old\dots» (función de riesgo), problema~8.1.
\newblock \url{https://ocw.mit.edu/courses/1-010-uncertainty-in-engineering-fall-2008/resources/app_08/}

\bibitem{mit16842}
O.~de Weck.
\newblock 16.842: Fundamentals of Systems Engineering.
\newblock Massachusetts Institute of Technology, MIT OpenCourseWare, otoño de 2015. Programa (modelo en V) y tareas A1 (partes interesadas, concepto de operación), A2 (requisitos y asignación de márgenes) y A4 (exploración de alternativas y selección del concepto).
\newblock \url{https://ocw.mit.edu/courses/16-842-fundamentals-of-systems-engineering-fall-2015/}

\bibitem{mit15401}
A.~Lo.
\newblock 15.401: Finance Theory I.
\newblock Massachusetts Institute of Technology, MIT OpenCourseWare, otoño de 2008. Sesiones «Present Value Relations» y «Capital Budgeting».
\newblock \url{https://ocw.mit.edu/courses/15-401-finance-theory-i-fall-2008/}

\bibitem{mit6172}
C.~Leiserson y J.~Shun.
\newblock 6.172: Performance Engineering of Software Systems.
\newblock Massachusetts Institute of Technology, MIT OpenCourseWare, otoño de 2018. Clase~10, «Measurement and Timing» (transcripción).
\newblock \url{https://ocw.mit.edu/courses/6-172-performance-engineering-of-software-systems-fall-2018/}

\bibitem{montgomery}
D.~C. Montgomery.
\newblock \emph{Design and Analysis of Experiments}, 9.ª ed.
\newblock Wiley, 2017. Cap.~2, «Simple Comparative Experiments».

\bibitem{alphaevolve}
A.~Novikov, N.~V{\~u}, M.~Eisenberger, E.~Dupont, P.-S. Huang, A.~Z. Wagner \emph{et al.}
\newblock AlphaEvolve: A coding agent for scientific and algorithmic discovery.
\newblock Prepublicación, Google DeepMind, 2025.
\newblock \url{https://arxiv.org/abs/2506.13131}

\bibitem{openevolve}
Desarrolladores de OpenEvolve.
\newblock OpenEvolve: implementación abierta inspirada en AlphaEvolve. Repositorio oficial; licencia Apache~2.0.
\newblock \url{https://github.com/algorithmicsuperintelligence/openevolve}

\bibitem{mipro}
K.~Opsahl-Ong, M.~J. Ryan, J.~Purtell, D.~Broman, C.~Potts, M.~Zaharia y O.~Khattab.
\newblock Optimizing instructions and demonstrations for multi-stage language model programs.
\newblock En \emph{Proceedings of EMNLP}, 2024.
\newblock \url{https://arxiv.org/abs/2406.11695}

\bibitem{mit1805}
J.~Orloff y J.~Bloom.
\newblock 18.05: Introduction to Probability and Statistics.
\newblock Massachusetts Institute of Technology, MIT OpenCourseWare, primavera de 2022. Lectura~24, «Bootstrap confidence intervals».
\newblock \url{https://ocw.mit.edu/courses/18-05-introduction-to-probability-and-statistics-spring-2022/}

\bibitem{enas}
H.~Pham, M.~Y. Guan, B.~Zoph, Q.~V. Le y J.~Dean.
\newblock Efficient neural architecture search via parameter sharing.
\newblock En \emph{Proceedings of the 35th International Conference on Machine Learning (ICML)}, 2018.
\newblock \url{https://arxiv.org/abs/1802.03268}

\bibitem{cs224u}
C.~Potts.
\newblock CS224U: Natural Language Understanding.
\newblock Stanford University, primavera de 2023. Tarea «Few-shot OpenQA with DSPy».
\newblock \url{https://web.stanford.edu/class/cs224u/}

\bibitem{ray}
P.~Moritz, R.~Nishihara, S.~Wang \emph{et al.}
\newblock Ray: A distributed framework for emerging AI applications.
\newblock En \emph{13th USENIX Symposium on Operating Systems Design and Implementation (OSDI)}, 2018. Documentación de Ray Tune y Ray Serve; licencia Apache~2.0.
\newblock \url{https://docs.ray.io/}

\bibitem{real}
E.~Real, A.~Aggarwal, Y.~Huang y Q.~V. Le.
\newblock Regularized evolution for image classifier architecture search.
\newblock En \emph{Proceedings of the AAAI Conference on Artificial Intelligence}, 2019.
\newblock \url{https://arxiv.org/abs/1802.01548}

\bibitem{aiact}
Parlamento Europeo y Consejo de la Unión Europea.
\newblock Reglamento (UE) 2024/1689, de 13 de junio de 2024, por el que se establecen normas armonizadas en materia de inteligencia artificial (Reglamento de Inteligencia Artificial).
\newblock \emph{Diario Oficial de la Unión Europea}, 12 de julio de 2024. Texto consultado, con sus modificaciones, en \url{https://artificialintelligenceact.eu/} (artículos~6, 14 y~113; anexos~I y~III).

\bibitem{ru}
R.~T. Rockafellar y S.~Uryasev.
\newblock Optimization of conditional value-at-risk.
\newblock \emph{Journal of Risk}, 2(3):21--41, 2000.

\bibitem{funsearch}
B.~Romera-Paredes, M.~Barekatain, A.~Novikov, M.~Balog, M.~P. Kumar, E.~Dupont, F.~J.~R. Ruiz, J.~S. Ellenberg, P.~Wang, O.~Fawzi, P.~Kohli y A.~Fawzi.
\newblock Mathematical discoveries from program search with large language models.
\newblock \emph{Nature}, 625:468--475, 2024.

\bibitem{sculley}
D.~Sculley, G.~Holt, D.~Golovin, E.~Davydov, T.~Phillips, D.~Ebner, V.~Chaudhary, M.~Young, J.-F. Crespo y D.~Dennison.
\newblock Hidden technical debt in machine learning systems.
\newblock En \emph{Advances in Neural Information Processing Systems 28 (NIPS)}, 2015.

\bibitem{cs329s}
C.~Huyen y equipo docente de CS329S.
\newblock CS329S: Machine Learning Systems Design.
\newblock Stanford University, invierno de 2022. Clases~1 «Understanding machine learning production», 2 «ML and Data Systems Fundamentals», 6 «Offline evaluation», 8 «Deployment» y~10 «Diagnosis of ML system failures, data distribution shifts \& monitoring», con notas de clase.
\newblock \url{https://stanford-cs329s.github.io/syllabus.html}

\bibitem{smith}
J.~E. Smith y R.~L. Winkler.
\newblock The optimizer's curse: Skepticism and postdecision surprise in decision analysis.
\newblock \emph{Management Science}, 52(3):311--322, 2006.

\bibitem{autoweka}
C.~Thornton, F.~Hutter, H.~H. Hoos y K.~Leyton-Brown.
\newblock Auto-WEKA: Combined selection and hyperparameter optimization of classification algorithms.
\newblock En \emph{Proceedings of the 19th ACM SIGKDD International Conference on Knowledge Discovery and Data Mining}, 2013.

\bibitem{wang-sc}
X.~Wang, J.~Wei, D.~Schuurmans, Q.~Le, E.~Chi, S.~Narang, A.~Chowdhery y D.~Zhou.
\newblock Self-consistency improves chain of thought reasoning in language models.
\newblock En \emph{International Conference on Learning Representations (ICLR)}, 2023.
\newblock \url{https://arxiv.org/abs/2203.11171}

\bibitem{ws}
D.~P. Williamson y D.~B. Shmoys.
\newblock \emph{The Design of Approximation Algorithms}.
\newblock Cambridge University Press, 2011. Edición electrónica libre. §3.3, «The bin-packing problem».
\newblock \url{https://www.designofapproxalgs.com/}

\bibitem{opro}
C.~Yang, X.~Wang, Y.~Lu, H.~Liu, Q.~V. Le, D.~Zhou y X.~Chen.
\newblock Large language models as optimizers.
\newblock Prepublicación, 2023 (versión~3, 2024).
\newblock \url{https://arxiv.org/abs/2309.03409}

\bibitem{d2l}
A.~Zhang, Z.~C. Lipton, M.~Li y A.~J. Smola.
\newblock \emph{Dive into Deep Learning}, versión 1.0.3.
\newblock Cambridge University Press, 2023. §19.1 «What Is Hyperparameter Optimization?» y §19.4 «Multi-Fidelity Hyperparameter Optimization».
\newblock \url{https://d2l.ai/}

\bibitem{zoph}
B.~Zoph y Q.~V. Le.
\newblock Neural architecture search with reinforcement learning.
\newblock En \emph{International Conference on Learning Representations (ICLR)}, 2017.
\newblock \url{https://arxiv.org/abs/1611.01578}

\end{thebibliography}


\end{document}
